% ============================================================
% FRAGMENTO -- continuación del capítulo "Diseño del software".
% NO es un documento standalone: no lleva \documentclass ni portada.
% Insértalo con % ============================================================
% FRAGMENTO -- continuación del capítulo "Diseño del software".
% NO es un documento standalone: no lleva \documentclass ni portada.
% Insértalo con % ============================================================
% FRAGMENTO -- continuación del capítulo "Diseño del software".
% NO es un documento standalone: no lleva \documentclass ni portada.
% Insértalo con % ============================================================
% FRAGMENTO -- continuación del capítulo "Diseño del software".
% NO es un documento standalone: no lleva \documentclass ni portada.
% Insértalo con \input{capitulo_diseno_software} (o copia/pega el
% contenido) dentro de tu documento original, en el punto donde va este
% capítulo. Ajusta el nivel de \section por \subsection si en tu
% documento esto ya cuelga de un \chapter{Diseño del software}.
%
% Requisitos que van en el PREAMBULO de tu documento maestro (no aquí):
%
%   \usepackage{amsmath,amssymb}
%   \usepackage{booktabs}
%   \usepackage{array}
%   \usepackage{xcolor}
%   \usepackage{listings}
%   \usepackage{float}
%   \usepackage{lmodern}    % fuentes escalables -- requisito de microtype
%                           % si tu documento usa babel en espanol (que
%                           % carga fuentes EC de mapa de bits sin esto)
%   \usepackage{microtype}  % reduce cajas "overfull" en parrafos con
%                           % identificadores \texttt largos (p.ej.
%                           % FUEL_CALIBRATION_FACTOR); combinar con:
%   \tolerance=1500
%   \emergencystretch=3em
%
%   \definecolor{codebg}{RGB}{245,245,245}
%   \definecolor{codekw}{RGB}{0,90,170}
%   \definecolor{codecom}{RGB}{110,110,110}
%   \definecolor{codestr}{RGB}{160,40,40}
%   \lstdefinestyle{cpp}{
%     language=C++, backgroundcolor=\color{codebg}, basicstyle=\ttfamily\small,
%     keywordstyle=\color{codekw}\bfseries, commentstyle=\color{codecom}\itshape,
%     stringstyle=\color{codestr}, numbers=left, numberstyle=\tiny\color{gray},
%     numbersep=8pt, breaklines=true, showstringspaces=false,
%     frame=single, rulecolor=\color{gray!40}, tabsize=2,
%   }
%   \lstdefinestyle{plain}{
%     backgroundcolor=\color{codebg}, basicstyle=\ttfamily\small,
%     breaklines=true, showstringspaces=false,
%     frame=single, rulecolor=\color{gray!40},
%   }
%   \lstset{style=cpp}
%
% Bibliografía: las fuentes citadas aquí (\cite{...}) están en
% referencias.bib, para fusionar con la bibliografía única de tu
% documento (donde ya tengas \bibliography{...} o \addbibresource{...}).
% \cite{} es el comando base de LaTeX, compatible con bibliografía clásica
% (\bibliographystyle{apalike} da formato "(Autor, Año)" sin paquetes
% extra), con natbib, y con biblatex.
% ============================================================

Esta sección retoma el capítulo de diseño de software y describe, parte por
parte, el firmware que corre en el ESP32-S3: las herramientas con las que se
construyó, cómo reparte el trabajo entre tareas, el protocolo con el que
habla con el vehículo, el método de cálculo del consumo, cómo persiste los
datos, cómo expone un panel de control, y el criterio detrás de cada
elección de componente y de librería.

% ============================================================
\section{Herramientas y entorno de desarrollo}
\label{sec:herramientas}
% ============================================================

\subsection{Visual Studio Code + PlatformIO}
El firmware se desarrolla en Visual Studio Code con la extensión PlatformIO
IDE, en vez del IDE oficial de Arduino. La diferencia no es cosmética:

\begin{itemize}
  \item \textbf{Toolchain embebido integrado.} PlatformIO gestiona la
        descarga y versión exacta del framework
        (\texttt{framework-arduinoespressif32}) \cite{arduinoesp32core2024},
        el compilador cruzado (\texttt{toolchain-xtensa-esp32s3}) y
        \texttt{esptool} para programar, todo declarado y fijado en
        \texttt{platformio.ini} \cite{platformio2024}. Esto hace que el
        \textit{build} sea reproducible entre máquinas —un requisito real
        una vez que el proyecto tiene múltiples archivos fuente y
        dependencias externas—, algo que el modelo de "sketch único" del
        IDE de Arduino maneja peor.
  \item \textbf{Autocompletado real de C++.} Con un servidor de lenguaje
        indexando todo el árbol de \texttt{src/}/\texttt{include/}, moverse
        entre los siete archivos fuente del proyecto y sus cabeceras
        compartidas (\texttt{config.h}, \texttt{types.h}) es inmediato.
  \item \textbf{Terminal integrado.} Compilar, subir y abrir el monitor
        serie desde el mismo terminal donde se edita el código evita el
        cambio de contexto entre editor y herramienta externa, y deja el
        registro de errores de compilación con enlaces directos a archivo y
        línea.
\end{itemize}

\subsection{Por qué FreeRTOS}
Un sistema que a la vez sondea un bus CAN sensible a tiempos, refresca una
pantalla, escribe en una tarjeta SD y sirve un panel web por WiFi es, por
naturaleza, un problema \textbf{concurrente}. Hay dos formas razonables de
resolverlo: un único \texttt{loop()} con máquinas de estado manuales y
sondeo no bloqueante en cada punto donde se podría esperar —viable en
proyectos pequeños, pero frágil a medida que crecen—, o un sistema operativo
de tiempo real con prioridades y primitivas de sincronización reales. Este
proyecto usa la segunda opción.

FreeRTOS no es, en rigor, una dependencia externa que se "agregó" —es el
planificador que el propio núcleo Arduino-ESP32 ya usa por debajo
\cite{freertos2024}; usarlo explícitamente es aprovechar la plataforma en
vez de pelear contra ella con un \texttt{loop()} monolítico encima.
Concretamente, esto permite:

\begin{itemize}
  \item Darle al sondeo CAN la \textbf{prioridad más alta} de todas las
        tareas, garantizando que una transacción I2C lenta hacia el OLED o
        una descarga de CSV por WiFi nunca retrase la lectura del vehículo.
  \item Proteger el estado compartido con mutex reales en vez de banderas
        improvisadas.
  \item Que agregar una funcionalidad nueva sea, casi siempre, agregar una
        tarea nueva con una prioridad clara, en vez de insertar más lógica
        condicional dentro de un bucle ya complejo —el modo de bajo consumo
        de la Sección~\ref{sec:idle} solo requirió tocar el bucle de la
        tarea CAN y agregar una bandera de estado que las demás tareas ya
        sabían consultar.
\end{itemize}

\begin{figure}[H]
\begin{lstlisting}[caption={Creación de tareas FreeRTOS en \texttt{src/main.cpp}.}, label={fig:tasks}]
statusLedsInit();
xTaskCreatePinnedToCore(statusLedsTask, "leds", STACK_LED_TASK, nullptr,
                         PRIO_LED_TASK, nullptr, CORE_LED_TASK);

bool oledOk = displayInit();
bool sdOk   = sdLoggerInit();
bool canOk  = canObd2Init();

webServerInit(); // reporta ErrorCode::WIFI_FAIL si el softAP falla

if (canOk) {
  xTaskCreatePinnedToCore(canObd2Task, "canObd2", STACK_CAN_TASK, nullptr,
                           PRIO_CAN_TASK, nullptr, CORE_CAN_TASK);
}
if (oledOk) {
  xTaskCreatePinnedToCore(displayTask, "display", STACK_DISPLAY_TASK, nullptr,
                           PRIO_DISPLAY_TASK, nullptr, CORE_DISPLAY_TASK);
}
if (sdOk) {
  xTaskCreatePinnedToCore(sdLoggerTask, "sdLogger", STACK_SD_TASK, nullptr,
                           PRIO_SD_TASK, nullptr, CORE_SD_TASK);
}
\end{lstlisting}
\end{figure}

% ============================================================
\section{Interfaz con el hardware}
% ============================================================

El software está organizado alrededor de las interfaces físicas del
ESP32-S3 \cite{espressif2024datasheet,espressif2024trm}. Cada periférico se
eligió por una razón concreta, no por ser "lo que venía en el kit":

\begin{table}[H]
\centering
\caption{Mapa de pines definido en \texttt{include/config.h}.}
\begin{tabular}{@{}lll@{}}
\toprule
Bus & Señal & GPIO ESP32-S3 \\
\midrule
TWAI (CAN)   & TX $\rightarrow$ SN65HVD230 CTX & GPIO4 \\
TWAI (CAN)   & RX $\leftarrow$ SN65HVD230 CRX  & GPIO5 \\
I2C          & SDA $\rightarrow$ OLED          & GPIO8 \\
I2C          & SCL $\rightarrow$ OLED          & GPIO9 \\
SPI (SD)     & SCK                              & GPIO12 \\
SPI (SD)     & MISO                             & GPIO13 \\
SPI (SD)     & MOSI                             & GPIO11 \\
SPI (SD)     & CS                               & GPIO10 \\
LED          & CAN activo                       & GPIO15 \\
LED          & Actividad SD                     & GPIO16 \\
LED          & Servidor activo                  & GPIO17 \\
LED          & Error                            & GPIO18 \\
\bottomrule
\end{tabular}
\end{table}

\textit{Nota.} Pines reservados —no usar—: GPIO 26-37 (bus interno hacia el
flash y la PSRAM octal del módulo N16R8) y GPIO 0/3/19/20/43-46 (pines de
\textit{strapping}, USB nativo y consola UART0 del chip).

\subsection{Controlador CAN: TWAI integrado, no un chip externo}
El ESP32-S3 trae un controlador CAN (TWAI) integrado en el propio chip
\cite{espressif2024trm}. Esa es la razón por la que el diseño no usa un
controlador CAN externo por SPI como el MCP2515, común en proyectos sobre
Arduino Uno/Nano que no tienen esta capacidad de fábrica: usar el
controlador nativo ahorra un dispositivo, una interrupción y un bus SPI que
ya está ocupado por la microSD, y evita la latencia adicional de tener que
serializar cada trama CAN a través de otro periférico intermediario.

\subsection{Transceptor CAN: SN65HVD230}
El TWAI del ESP32-S3 solo genera niveles lógicos TTL; convertirlos a las
señales diferenciales CAN\_H/CAN\_L del bus físico requiere un transceptor.
Se eligió el SN65HVD230 \cite{sn65hvd230_datasheet} en concreto porque opera
nativamente a 3.3\,V —el mismo nivel lógico del ESP32-S3—, evitando un
\textit{level-shifter} que sí haría falta con transceptores clásicos de 5\,V
como el MCP2551. CAN\_H y CAN\_L se conectan a los pines 6 y 14 del conector
OBD2 (J1962).

\subsection{Pantalla: SSD1306 OLED por I2C}
Se eligió un OLED monocromático de $128\times64$ px con controlador SSD1306
\cite{solomon2008ssd1306} sobre I2C en vez de, por ejemplo, una pantalla TFT
a color por SPI, por tres razones: consume solo 2 pines (I2C) dejando el SPI
libre para la SD; su consumo (10-20\,mA encendida, y se puede apagar por
software, Sección~\ref{sec:idle}) es bajo comparado con un TFT retroiluminado;
y la resolución alcanza de sobra para los pocos números grandes que hace
falta mostrar al conductor —no se necesita color ni gráficos complejos en el
tablero físico, eso ya lo cubre el panel web.

\subsection{Almacenamiento: microSD por SPI dedicado}
La tarjeta microSD usa su propio bus SPI, sin compartirlo con ningún otro
periférico. La alternativa —reutilizar un bus SPI ya en uso— se descartó
porque el registro en la SD ocurre una vez por segundo de forma predecible
mientras que una eventual escritura de mayor tamaño (por ejemplo, servir una
descarga de CSV completa) no debe competir por el bus con otra
responsabilidad no relacionada; un bus dedicado simplifica el análisis de
tiempos y evita tener que coordinar dos periféricos distintos sobre el mismo
SPI.

% ============================================================
\section{Arquitectura de tareas}
% ============================================================

\begin{table}[H]
\centering
\caption{Tareas FreeRTOS, definidas en \texttt{include/config.h} y creadas en \texttt{src/main.cpp}.}
\begin{tabular}{@{}llccc@{}}
\toprule
Tarea & Función & Núcleo & Prioridad & Pila \\
\midrule
CAN/OBD2          & \texttt{canObd2Task}       & 0 & 5 & 4096 \\
SD logger         & \texttt{sdLoggerTask}      & 1 & 4 & 6144 \\
WebSocket push     & \texttt{webSocketPushTask} & 0 & 3 & 4096 \\
OLED display      & \texttt{displayTask}       & 1 & 2 & 4096 \\
LEDs de estado    & \texttt{statusLedsTask}    & 1 & 1 & 2048 \\
\bottomrule
\end{tabular}
\end{table}

El núcleo 0 concentra lo sensible a tiempos (sondeo CAN de alta prioridad +
la pila WiFi/AsyncTCP, que en Arduino-ESP32 corre en ese mismo núcleo). El
núcleo 1 queda con las tareas de interfaz/almacenamiento, que toleran más
\textit{jitter} sin afectar la lectura del vehículo. Las prioridades no son
arbitrarias: reflejan directamente qué tan caro es que cada tarea llegue
tarde (perder un ciclo de sondeo CAN corrompe el cálculo de consumo; llegar
tarde a refrescar un LED es imperceptible).

\subsection{Sincronización entre tareas}
Dos mutex protegen el estado compartido: \texttt{g\_dataMutex} protege
\texttt{g\_vehicleData} (la tarea CAN arma cada ciclo en una copia local y
solo toca la estructura global, bajo mutex, en el paso final de fusión —
Figura~\ref{fig:merge}); \texttt{g\_sdMutex} protege el bus SPI de la
microSD, compartido entre el \textit{logger} periódico y los endpoints de
descarga/gráfica del servidor web. El criterio para usar mutex en vez de,
por ejemplo, deshabilitar interrupciones o copiar datos sin protección, es
que ambos recursos (la estructura de datos y el bus SPI) tienen escritores y
lectores en tareas de distinta prioridad corriendo en núcleos distintos —un
mutex es la primitiva correcta quando el acceso puede, en principio,
bloquear brevemente sin romper los tiempos críticos del sistema.

\begin{figure}[H]
\begin{lstlisting}[caption={Fusión bajo mutex al final de cada ciclo, \texttt{src/can\_obd2.cpp}.}, label={fig:merge}]
if (xSemaphoreTake(g_dataMutex, pdMS_TO_TICKS(100)) == pdTRUE) {
  VehicleData &g = g_vehicleData;
  g.rpm = scratch.rpm;
  g.speed_kmh = scratch.speed_kmh;
  // ... resto de los campos ...
  g.can_status = status;
  g.data_valid = criticalOk;
  g.last_update_ms = now;

  if (criticalOk && dt > 0.0f && dt < 5.0f) {
    fuelCalcUpdate(g, dt);
  }
  xSemaphoreGive(g_dataMutex);
}
\end{lstlisting}
\end{figure}

\subsection{Manejo de errores}
\texttt{systemReportError(ErrorCode)} centraliza el reporte de fallos: guarda
el último código en \texttt{g\_systemStatus.error\_code} (dispara el parpadeo
del LED de error) y lo imprime por \texttt{Serial} con el detalle completo.
El criterio es separar dos audiencias distintas: el LED es para el
conductor, que solo necesita saber \textit{que algo falló}; el log serie es
para depuración en banco, donde sí importa el detalle completo y el
historial.

\begin{table}[H]
\centering
\caption{Códigos de \texttt{ErrorCode} (\texttt{include/types.h}).}
\begin{tabular}{@{}cl@{}}
\toprule
Código & Significado \\
\midrule
1 & Fallo al iniciar el driver TWAI/CAN \\
2 & Fallo al iniciar la microSD \\
3 & Fallo al iniciar el OLED \\
4 & Fallo al iniciar el punto de acceso WiFi \\
\bottomrule
\end{tabular}
\end{table}

% ============================================================
\section{Protocolo CAN / OBD2}
% ============================================================

El controlador TWAI del ESP32-S3 se configura a \textbf{500\,kbit/s}, el
estándar de la mayoría de vehículos livianos con OBD2 sobre CAN
(ISO 15765-4, \cite{iso157654}). Todas las
peticiones usadas caben en un único \textit{frame} CAN de 8 bytes
(\textit{single frame}), así que no hace falta implementar el mecanismo
multi-\textit{frame} de control de flujo ISO-TP —una simplificación
deliberada: ninguno de los PIDs que este proyecto necesita supera ese
tamaño, y soportar multi-\textit{frame} sin necesitarlo solo agrega
superficie de error.

\subsection{PIDs consultados (Modo 01)}
El conjunto de PIDs se eligió por ser el mínimo necesario para el cálculo
speed-density (Sección~\ref{sec:speeddensity}) más los que enriquecen el
registro sin costo relevante de tiempo de bus, siguiendo el estándar SAE
J1979 \cite{sae_j1979_2014}:

\begin{table}[H]
\centering
\caption{PIDs Mode 01 consultados por \texttt{canObd2Task}.}
\begin{tabular}{@{}cll@{}}
\toprule
PID & Dato & Uso \\
\midrule
\texttt{0x0C} & RPM                          & speed-density (crítico) \\
\texttt{0x0B} & MAP (kPa)                    & speed-density (crítico) \\
\texttt{0x0F} & IAT (\textdegree C)          & speed-density (crítico) \\
\texttt{0x0D} & Velocidad (km/h)             & L/100km, distancia (crítico) \\
\texttt{0x05} & Temperatura refrigerante     & mostrado en pantalla/CSV \\
\texttt{0x04} & Carga calculada (\%)         & mostrado en pantalla/CSV \\
\texttt{0x06} & Fuel trim corto plazo (\%)   & corrige el AFR estimado \\
\texttt{0x07} & Fuel trim largo plazo (\%)   & corrige el AFR estimado \\
\bottomrule
\end{tabular}
\end{table}

\textit{Nota.} Los cuatro PIDs marcados "crítico" son indispensables para el
cálculo speed-density; si la ECU no responde alguno, el ciclo completo se
marca inválido y no se integra al viaje.

\subsection{Formato de la petición}

\begin{figure}[H]
\begin{lstlisting}[style=plain,caption={Trama de petición Mode 01, PID genérico.}]
ID = 0x7DF
Data[0] = 0x02   ; SF, 2 bytes de payload (modo + pid)
Data[1] = 0x01   ; Modo 01: datos actuales
Data[2] = <PID>
Data[3..7] = 0x00 (relleno)
\end{lstlisting}
\end{figure}

Las ECUs responden en el rango \texttt{0x7E8}-\texttt{0x7EF}; se acepta la
primera respuesta cuyo \texttt{Data[1] == 0x41} y \texttt{Data[2]} coincida
con el PID pedido.

\subsection{Recuperación de bus-off}
Si el controlador entra en estado \texttt{BUS\_OFF}, la recuperación se
sondea en cada vuelta del bucle sin bloquear con una espera fija —según el
driver TWAI de ESP-IDF \cite{espressif2024trm}, la recuperación solo llega a
\texttt{STOPPED} tras contar 128 ocurrencias de 11 bits recesivos
consecutivos en el bus, un tiempo que depende del tráfico real y no se puede
predecir de antemano; sondear en vez de esperar un tiempo fijo evita tanto
reintentar demasiado pronto (fallar en vano) como esperar de más
innecesariamente.

% ============================================================
\section{Método de cálculo: Speed-Density}
\label{sec:speeddensity}
% ============================================================

El método speed-density estima el flujo másico de aire admitido por el motor
a partir de la ley de gases ideales, sin medirlo con un sensor MAF —la
razón de fondo para elegir este método sobre uno basado en MAF real es que
la mayoría de los vehículos con OBD2 reportan MAP, no flujo de aire directo,
así que speed-density es el único método viable sin instalar un sensor
adicional—:

\begin{equation}
\dot{m}_{aire} \;[\text{g/s}] =
\frac{\dfrac{VE}{100} \cdot RPM \cdot MAP\,[\text{kPa}] \cdot V_d\,[\text{L}] \cdot 28.97}
     {2 \cdot 8.314 \cdot T_{IAT}\,[\text{K}] \cdot 60}
\end{equation}

donde $VE$ es la eficiencia volumétrica (\%) interpolada de una tabla
calibrable (Sección~\ref{sec:ve}); $V_d$ la cilindrada en litros; $28.97$ la
masa molar del aire (g/mol); $8.314$ la constante universal de los gases
(L$\cdot$kPa/(mol$\cdot$K)); el factor $2$ corresponde a que un motor de 4
tiempos completa una admisión cada 2 vueltas de cigüeñal; y $T_{IAT}$ la
temperatura de admisión en Kelvin, con guarda de rango
[$-40$,\,$100$]\,\textdegree C.

\begin{figure}[H]
\begin{lstlisting}[caption={Núcleo del cálculo speed-density, \texttt{src/fuel\_calc.cpp}.}, label={fig:fuelcalc}]
float ve = fuelCalcVeForRpm(vd.rpm);
float iat_K = vd.iat_c + 273.15f;
if (iat_K < 233.15f || iat_K > 373.15f) iat_K = 288.15f;

const float MM_AIR = 28.97f; // g/mol
const float R = 8.314f;      // L*kPa/(mol*K)

float maf = (ve / 100.0f) * (float)vd.rpm * vd.map_kpa *
            ENGINE_DISPLACEMENT_L * MM_AIR /
            (2.0f * R * iat_K * 60.0f);
\end{lstlisting}
\end{figure}

\subsection{De aire a combustible}

\begin{equation}
\dot{m}_{comb}\,[\text{g/s}] = \frac{\dot{m}_{aire}}{AFR} \cdot k_{cal}
\end{equation}

$AFR$ es la relación aire/combustible, por defecto el valor estequiométrico
de la gasolina ($AFR_{stoich} = 14.7$), opcionalmente corregido con los
\textit{fuel trims} que reporta la ECU:

\begin{equation}
AFR_{efectivo} = \frac{AFR_{stoich}}{1 + \frac{STFT\% + LTFT\%}{100}}
\end{equation}

acotado a $[8, 25]$. $k_{cal}$ es \texttt{FUEL\_CALIBRATION\_FACTOR}
(Sección~\ref{sec:calibracion}), $1.0$ por defecto. El criterio para incluir
la corrección por \textit{fuel trim} es que, sin ella, el cálculo asume que
el motor siempre corre exactamente en estequiométrico, cosa que en la
práctica no ocurre; usar los \textit{trims} que la propia ECU ya reporta es
una corrección "gratis" —no requiere sensores adicionales— frente a un error
sistemático conocido.

\subsection{De masa a volumen}

\begin{align}
\dot{V}_{comb}\,[\text{L/h}] &= \frac{\dot{m}_{comb} \cdot 3600}{\rho_{comb}} \\[4pt]
L/100km &=
  \begin{cases}
    \dfrac{\dot{V}_{comb}}{v}\cdot 100 & \text{si } v > 2\,\text{km/h} \\[6pt]
    0 & \text{vehículo detenido}
  \end{cases}
\end{align}

con $\rho_{comb} = 745\,\text{g/L}$ y $v$ la velocidad en km/h.

\subsection{Integración del viaje}

\begin{align}
L_{viaje} &\mathrel{+}= \dot{m}_{comb} \cdot dt \,/\, \rho_{comb} \\
km_{viaje} &\mathrel{+}= v \cdot dt / 3600 \\
Bs_{viaje} &= L_{viaje} \cdot P_{L}
\end{align}

con $dt$ el paso de tiempo real entre ciclos —no un valor fijo, precisamente
para no perder precisión si un ciclo de sondeo tarda más o menos de lo
típico— y $P_L = 6.96$ (\texttt{FUEL\_PRICE\_BS\_PER\_L}) el precio de
referencia del litro.

\subsection{Tabla de eficiencia volumétrica (VE)}
\label{sec:ve}

\begin{table}[H]
\centering
\caption{\texttt{kVeTable}, \texttt{src/fuel\_calc.cpp}. Interpolación lineal entre puntos.}
\begin{tabular}{@{}cc@{}}
\toprule
RPM & VE (\%) \\
\midrule
700  & 60.0 \\
1500 & 74.0 \\
2500 & 85.0 \\
4000 & 88.0 \\
5500 & 83.0 \\
7000 & 74.0 \\
\bottomrule
\end{tabular}
\end{table}

\textit{Nota.} Esta tabla es un punto de partida genérico, no un dato de
fábrica del vehículo. Sin calibrarla, el consumo calculado tendrá un error
sistemático.

\subsection{Calibración contra un tanque real}
\label{sec:calibracion}
No existe forma de que el ESP32 sepa, por sí solo, si el número que calcula
es correcto —no hay con qué comparar dentro del propio microcontrolador—.
El procedimiento recomendado: llenar el tanque y anotar el odómetro, manejar
con normalidad hasta la próxima recarga, comparar los litros reales cargados
contra \texttt{trip\_fuel\_L} acumulado en esa misma distancia, y ajustar

\begin{equation}
\texttt{FUEL\_CALIBRATION\_FACTOR} = \frac{L_{real}}{L_{calculado}}
\end{equation}

en \texttt{config.h}. El criterio de aplicar un único factor multiplicativo
sobre \texttt{fuel\_gs}, en vez de reajustar toda la tabla VE a mano, es que
corrige de forma consistente el instantáneo, el acumulado del viaje y el
costo con un solo número, y es reversible/auditable —a diferencia de
retocar seis puntos de la tabla VE por prueba y error.

% ============================================================
\section{Registro en microSD}
% ============================================================

Cada arranque crea un archivo nuevo \texttt{/logs/trip\_NNN.csv} (índice
autoincremental calculado escaneando los archivos existentes, en vez de un
contador separado que podría desincronizarse tras un corte de energía). Se
escribe una fila por segundo, solo después del primer ciclo CAN válido, y se
pausa por completo durante el modo de bajo consumo (Sección~\ref{sec:idle}).

\subsection{Por qué CSV}
Se eligió CSV en vez de un formato binario compacto o JSON por fila por tres
razones: es directamente legible/abrible en cualquier hoja de cálculo sin
escribir un parser; el tamaño de fila (unos 100 bytes) es intrascendente
frente a la capacidad de una microSD moderna incluso a una fila por segundo
durante horas; y es el formato que el propio panel web ya necesita poder
descargar y volver a interpretar (Sección~\ref{sec:panelweb}), así que un
mismo formato sirve para ambos usos sin conversión intermedia.

\begin{figure}[H]
\begin{lstlisting}[style=plain,caption={Fila de ejemplo del CSV (encabezado + una fila de datos).}]
uptime_ms,rpm,speed_kmh,map_kpa,iat_c,coolant_c,ve_pct,maf_gs,
fuel_gs,instant_Lh,instant_L100km,trip_km,trip_fuel_L,
trip_avg_L100km,can_status,trip_cost_bs

184532,2150,62.0,58.0,29.0,89.0,81.4,7.82,0.532,1.91,3.09,
12.400,0.384,3.10,1,2.67
\end{lstlisting}
\end{figure}

\begin{table}[H]
\centering
\caption{Columnas escritas por \texttt{sdLoggerTask} (\texttt{src/sd\_logger.cpp}).}
\begin{tabular}{@{}ll@{}}
\toprule
Columna & Descripción \\
\midrule
\texttt{uptime\_ms}        & Milisegundos desde el arranque \\
\texttt{rpm}                & Revoluciones por minuto \\
\texttt{speed\_kmh}        & Velocidad, km/h \\
\texttt{map\_kpa}          & Presión absoluta del múltiple, kPa \\
\texttt{iat\_c}             & Temperatura del aire de admisión, \textdegree C \\
\texttt{coolant\_c}        & Temperatura del refrigerante, \textdegree C \\
\texttt{ve\_pct}            & Eficiencia volumétrica usada, \% \\
\texttt{maf\_gs}            & Flujo másico de aire calculado, g/s \\
\texttt{fuel\_gs}           & Flujo másico de combustible calculado, g/s \\
\texttt{instant\_Lh}       & Consumo instantáneo, L/h \\
\texttt{instant\_L100km}  & Consumo instantáneo, L/100km \\
\texttt{trip\_km}          & Distancia acumulada del viaje, km \\
\texttt{trip\_fuel\_L}    & Combustible acumulado del viaje, L \\
\texttt{trip\_avg\_L100km} & Promedio del viaje, L/100km \\
\texttt{can\_status}       & Estado del bus CAN (enum \texttt{CanStatus}) \\
\texttt{trip\_cost\_bs}    & Costo acumulado del viaje, Bs \\
\bottomrule
\end{tabular}
\end{table}

% ============================================================
\section{Modo de bajo consumo ("sin contacto")}
\label{sec:idle}
% ============================================================

\subsection{Motivación}
El dispositivo se alimenta del pin 16 del conector OBD2, que entrega
12\,V \textbf{constantes} sin importar la posición de la llave. Si el
firmware siguiera sondeando el bus a ritmo completo, con el panel web y
todos los periféricos encendidos las 24 horas, se convertiría en una carga
parásita permanente sobre la batería del vehículo. La primera versión de
este modo se basaba en detectar RPM$=0$ y solo bajaba el ritmo de sondeo
—una mejora real pero modesta, porque el ESP32-S3 seguía activo con el
WiFi encendido—. La versión actual va más lejos: separa la alimentación en
dos rieles y usa \textbf{deep sleep} real del ESP32-S3 mientras no hay
contacto.

\subsection{Dos rieles de alimentación}
La PCB entrega dos salidas de 3.3\,V independientes desde el 12\,V
constante del OBD2:

\begin{itemize}
  \item \textbf{Riel 1} (siempre encendido): ESP32-S3 + SN65HVD230. Tiene
        que quedar siempre alimentado porque el transceptor es lo único
        capaz de "escuchar" el bus CAN para saber cuándo volvió el
        contacto.
  \item \textbf{Riel 2} (conmutado por software): LEDs + OLED, y por
        separado la microSD. Se corta por completo mientras no hay
        contacto.
\end{itemize}

\subsection{Circuito de \textit{power-gating}: MOSFET y pines}
\label{sec:powergate}
Cada riel 2 se controla con un MOSFET de \textbf{canal P, como interruptor
de alto lado} (corta el 3.3\,V que entra al periférico, no su GND). Se
prefiere sobre un N-MOSFET de bajo lado porque con este último el GND del
periférico puede quedar ligeramente flotante mientras conduce, lo que
genera problemas de integridad de señal en buses referenciados a GND como
I2C (OLED) y SPI (microSD).

\begin{table}[H]
\centering
\caption{Pines de \textit{power-gating}, \texttt{include/config.h}.}
\begin{tabular}{@{}lll@{}}
\toprule
Señal & GPIO ESP32-S3 & Controla \\
\midrule
\texttt{PIN\_PWR\_LED\_OLED} & GPIO6 & LEDs + OLED \\
\texttt{PIN\_PWR\_SD}        & GPIO7 & microSD \\
\bottomrule
\end{tabular}
\end{table}

Ambos GPIO son de dominio RTC en el ESP32-S3 (rango GPIO0-21), requisito
para que \texttt{gpio\_hold\_en()} pueda retener su nivel lógico durante el
\textit{deep sleep} \cite{espressif2024trm}. Topología de cada MOSFET
(misma para los dos rieles):

\begin{figure}[H]
\begin{lstlisting}[style=plain,caption={Interruptor de alto lado con P-MOSFET, un canal por riel.}]
Riel 2 (3.3V) ---+--------------------- Source (P-MOSFET)
                 |
              [R_pullup ~10k]
                 |
GPIO --[R_gate ~1k]------------------- Gate
                                          |
                                        Drain --- carga (OLED+LED, o microSD)
\end{lstlisting}
\end{figure}

El criterio detrás de cada componente: como el GPIO y el riel 2 comparten
el mismo nivel lógico (3.3\,V), no hace falta un transistor adicional para
manejar la compuerta —GPIO en alto da $V_{gs}=0$ (MOSFET apagado); GPIO en
bajo da $V_{gs}\approx-3.3\,\text{V}$, muy por debajo del umbral típico de
un P-MOSFET de nivel lógico ($-1$ a $-2\,\text{V}$ aprox., p.\,ej. un
AO3401A o equivalente)—. La resistencia de \textit{pull-up} en la
compuerta asegura que el MOSFET arranque apagado por defecto durante el
breve instante antes de que el firmware configure el GPIO como salida; la
resistencia en serie hacia el GPIO limita la corriente de carga de la
compuerta.

\subsection{Cadena de regulación}
\label{sec:regchain}
La alimentación completa, del 12\,V del OBD2 a los dos rieles de 3.3\,V,
pasa por tres reguladores en cascada:

\begin{enumerate}
  \item \textbf{TPS54302}: convertidor \textit{buck} (conmutado) que baja
        el 12\,V del pin 16 a 5\,V. Es la única etapa de conversión
        conmutada del sistema; las dos siguientes son lineales.
  \item \textbf{AMS1117} (5\,V $\rightarrow$ 3.3\,V): alimenta el
        \textbf{riel 1} (ESP32-S3 + SN65HVD230). Nunca se apaga.
  \item \textbf{AP2112} (5\,V $\rightarrow$ 3.3\,V): alimenta el
        \textbf{riel 2}; su salida de 3.3\,V entra a los dos MOSFET de la
        Sección~\ref{sec:powergate}, que gatean OLED+LED y microSD por
        separado. El AP2112 en sí mismo nunca se apaga —solo se corta lo
        que hay \textit{después} de los MOSFET—, pero al ser una LDO de
        muy baja corriente de reposo su aporte al consumo en espera es
        pequeño de todos modos (Sección~\ref{sec:consumo}).
\end{enumerate}

Esta cadena (\textit{buck} + 2 LDO en paralelo desde el mismo 5\,V) es una
elección de diseño válida y común, pero tiene una consecuencia directa para
el consumo "sin contacto" que conviene tener presente: a diferencia del
AP2112, el \textbf{AMS1117 es una LDO de propósito general, no de bajo
consumo} —su corriente de reposo (\textit{ground current}) es
comparativamente alta, típicamente 5-8\,mA casi sin importar la carga en
su salida, con variación notable entre fabricantes clon del mismo chip—.
Como el AMS1117 alimenta justo el riel que nunca se apaga, esa corriente de
reposo se sostiene las 24\,horas del día, se detalla en la
Sección~\ref{sec:consumo}.

\subsection{Implementación en firmware}
\texttt{canObd2CheckContact()} manda un único \textit{ping} (PID
\texttt{0x0C}, RPM, obligatorio en todo ECU OBD2) y reporta si algo
respondió, sin importar el valor. Se usa en dos puntos:

\begin{itemize}
  \item Al arrancar (\texttt{setup()}, tanto en un encendido normal como al
        despertar del \textit{deep sleep}): si no hay contacto, el chip
        vuelve a dormir antes de tocar el riel 2 o el WiFi.
  \item Ya despierto, dentro de \texttt{canObd2Task}: si el bus deja de
        responder durante \texttt{CONTACT\_LOST\_DEBOUNCE\_CYCLES} ciclos
        seguidos (3, para no dormirse por un solo \textit{frame} perdido en
        un bus ocupado), se considera que se fue el contacto.
\end{itemize}

\begin{figure}[H]
\begin{lstlisting}[caption={Entrada a deep sleep, \texttt{src/power\_sleep.cpp}.}, label={fig:deepsleep}]
void enterDeepSleepUntilContact() {
  sdLoggerShutdown();  // cierra el archivo antes de cortarle la energia
  WiFi.mode(WIFI_OFF);

  powerGateSet(false);
  gpio_hold_en((gpio_num_t)PIN_PWR_LED_OLED);
  gpio_hold_en((gpio_num_t)PIN_PWR_SD);
  gpio_deep_sleep_hold_en(); // sin esto el hold no sobrevive al deep sleep

  esp_sleep_enable_timer_wakeup(DEEP_SLEEP_WAKE_INTERVAL_US);
  esp_deep_sleep_start();  // no retorna
}
\end{lstlisting}
\end{figure}

Cerrar el archivo de la SD \textbf{antes} de cortar su MOSFET es
obligatorio, no una formalidad: cortar la alimentación con un archivo
abierto arriesga corromper el sistema de archivos completo, no solo perder
la última fila. El \textit{deep sleep} borra toda la RAM salvo el dominio
RTC, así que al despertar \texttt{setup()} vuelve a correr desde el
principio, exactamente como un arranque en frío.

\subsection{Consumo estimado y autonomía de la batería}
\label{sec:consumo}
Los siguientes números son estimaciones a partir de valores típicos de
hoja de datos, no mediciones sobre la placa real —sirven para tener una
idea de orden de magnitud, no un valor exacto—.

\begin{table}[H]
\centering
\small
\caption{Contribuyentes al consumo promedio "sin contacto", cadena real (TPS54302 + AMS1117 + AP2112).}
\begin{tabular}{@{}p{1.7cm}p{3.6cm}p{1.7cm}p{4.8cm}@{}}
\toprule
Nodo & Fuente & Aporte típico & Nota \\
\midrule
3.3\,V (riel 1) & ESP32-S3 en deep sleep & $\sim$0.02\,mA & 10-25\,$\mu$A \cite{espressif2024datasheet} \\
3.3\,V (riel 1) & SN65HVD230, siempre activo & $\sim$5-10\,mA & Modo normal, no standby \cite{sn65hvd230_datasheet} \\
3.3\,V (riel 1) & Pulsos de despertar (cada 20\,s) & $\sim$1.5-2\,mA & Promediado: $\sim$150\,mA por $\sim$0.2-0.3\,s \\
5\,V & \textbf{AMS1117}, corriente de reposo propia & \textbf{$\sim$5-8\,mA} & Casi independiente de la carga \\
5\,V & AP2112, corriente de reposo propia & $\sim$0.06\,mA & Sin carga (MOSFET de salida abiertos) \\
12\,V & TPS54302, corriente de reposo propia & $\sim$0.1-0.3\,mA & \\
\midrule
\multicolumn{2}{@{}l}{\textbf{Total aproximado, lado 12\,V}} & \textbf{$\sim$7-9\,mA} & Ver conversión abajo \\
\bottomrule
\end{tabular}
\end{table}

El resultado es contraintuitivo pero importante: \textbf{el transceptor CAN
y la LDO del riel 1, no el ESP32-S3, dominan el consumo mientras el
dispositivo "duerme"} —porque ese riel no se puede cortar sin perder la
capacidad de detectar cuándo vuelve el contacto (Sección~\ref{sec:powergate}).
El \textit{deep sleep} resuelve la parte del ESP32-S3 casi por completo,
pero no puede resolver la del transceptor ni la de su LDO con esta
arquitectura de dos rieles fija.

Convertir la corriente del lado 3.3\,V/5\,V al lado 12\,V: el AMS1117 y el
AP2112 son LDO (no reducen corriente, solo voltaje, disipando la diferencia
como calor), así que la corriente que el TPS54302 debe entregar en su
salida de 5\,V es aproximadamente la suma de lo que consumen el riel 1
completo ($\sim$8.5\,mA de ESP32+SN65+pulsos) más la corriente de reposo
propia del AMS1117 ($\sim$6.5\,mA punto medio) y del AP2112
($\sim$0.06\,mA): $\approx$15\,mA a 5\,V. El TPS54302, al ser un
\textit{buck} conmutado, sí convierte potencia de forma más eficiente
($\sim$80\,\% en carga liviana, estimado):

\begin{equation}
I_{12V} \approx \frac{5\,\text{V} \times 15\,\text{mA}}{12\,\text{V} \times 0.80} + 0.2\,\text{mA (I}_q\text{ TPS54302)} \approx 8\,\text{mA}
\end{equation}

Tomando esos $\sim$8\,mA y una batería típica de auto de 50\,Ah (de la cual,
por buena práctica, no conviene usar más del 50\,\% para no acortar su vida
útil, es decir $\sim$25\,000\,mAh utilizables):

\begin{equation}
\frac{25\,000\,\text{mAh}}{8\,\text{mA}} \approx 3\,125\,\text{horas} \approx 130\,\text{días}
\end{equation}

Es decir, \textbf{el dispositivo por sí solo} tardaría del orden de 4-5
meses en consumir la mitad de una batería sana —menos que si el riel 1
usara una LDO de baja corriente de reposo en vez del AMS1117, pero todavía
un margen amplio frente al uso normal del vehículo—. En la práctica esto
importa menos de lo que parece, porque el propio vehículo ya tiene su
propio consumo parásito de fábrica (alarma, memorias de radio/ECU, reloj,
etc.), típicamente $\sim$20-50\,mA con todo "apagado" —del mismo orden o
mayor que el aporte de este dispositivo—, y ese consumo de base es el que
realmente determina cuánto aguanta la batería estacionada, con o sin el
ESP32-S3 conectado.

\textbf{Decisión de diseño:} se evaluó reemplazar el AMS1117 del riel 1 por
una LDO de baja corriente de reposo (el propio AP2112, un MCP1700, o una
TPS7A02 con $I_q$ en el orden de nanoamperios bajarían el consumo "sin
contacto" varios mA más), pero se optó por mantener el AMS1117 y el AP2112
tal como estaban ya elegidos en la placa —es el contribuyente individual
más grande después del SN65HVD230, y queda documentado aquí como una
compensación consciente entre simplicidad/costo de la lista de materiales y
la autonomía teórica de la batería, no un descuido.

\subsection{Recomendación de mantenimiento}
Con esos números, la recomendación práctica no cambia mucho respecto al
cuidado normal de cualquier batería de auto: \textbf{arrancar/circular el
vehículo al menos cada 2-3 semanas}, o usar un cargador de mantenimiento
(\textit{battery tender}) si va a quedar estacionado más tiempo que eso.
Con la cadena de reguladores actual (AMS1117 incluido) el dispositivo
agrega una fracción del consumo parásito que el auto ya tiene de fábrica
—notable, pero no dominante—, así que este intervalo sigue siendo una
recomendación razonable sin necesitar ser más estricta.

\subsection{¿Afecta consultar tan seguido?}
Con la arquitectura actual hay dos frecuencias distintas que vale la pena
separar:

\begin{itemize}
  \item \textbf{Sondeo CAN mientras hay contacto} (un ciclo completo de 8
        PIDs cada $\sim$0.2-1.4\,s): el bus CAN a 500\,kbit/s tiene
        capacidad de sobra para esto —cada trama de petición/respuesta son
        unos 100 bits, así que incluso varias por segundo representan bien
        menos del 1\,\% del ancho de banda del bus—. Una ECU real está
        diseñada para atender consultas de diagnóstico bastante más
        frecuentes que esto sin ningún problema; no hay riesgo de saturar
        el bus ni de sobrecargar la ECU.
  \item \textbf{Escritura a la microSD, una fila por segundo}
        (\texttt{SD\_LOG\_INTERVAL\_MS}): las escrituras son secuenciales
        —se van agregando al final del archivo, no se reescribe una y otra
        vez el mismo bloque—, que es el mejor caso posible para el
        desgaste de la memoria flash. Una fila ronda los 100 bytes; a una
        por segundo son unos 360\,KB por hora de viaje, insignificante
        frente a la resistencia típica de una microSD moderna (miles de
        ciclos de escritura por bloque, con \textit{wear leveling}
        integrado). No es un factor de desgaste relevante en la práctica.
\end{itemize}

% ============================================================
\section{Panel web}
\label{sec:panelweb}
% ============================================================

\subsection{Punto de acceso WiFi: por qué SoftAP en vez de solo STA}
El dispositivo siempre levanta su propio punto de acceso (SSID
\texttt{ESP32\_OBD2}, IP típica \texttt{192.168.4.1}) en vez de depender
únicamente de unirse a una red WiFi conocida. El criterio es el contexto de
uso real: dentro de un vehículo casi nunca hay una red WiFi conocida en
rango, así que un diseño solo-STA dejaría el panel inaccesible la mayor
parte del tiempo. El SoftAP garantiza acceso local sin esa dependencia;
unirse además como cliente a una red conocida queda como algo opcional,
únicamente para sincronizar hora por NTP cuando hay una disponible.

\subsection{Por qué ESPAsyncWebServer en vez del servidor síncrono de Arduino}
Se eligió ESPAsyncWebServer \cite{espasyncwebserver2024} sobre el
\texttt{WebServer} síncrono incluido en el core de Arduino-ESP32 porque
atiende las conexiones de forma asíncrona, basada en eventos: una petición
lenta (por ejemplo, una descarga de CSV grande) no bloquea el hilo mientras
espera que el cliente reciba los datos. Dado que el servidor corre en el
mismo núcleo que la pila WiFi (Sección~\ref{sec:herramientas}), un modelo
síncrono habría significado que una descarga lenta retrasara también el
resto del tráfico de red, incluyendo el WebSocket que empuja los datos en
vivo.

\subsection{Empuje en vivo: WebSocket con respaldo por \textit{polling}}
El estado se empuja a los clientes por WebSocket cada 500\,ms
—Figura~\ref{fig:merge} ya mostró de dónde sale ese dato—, con un
mecanismo de respaldo: si el navegador detecta que el WebSocket no está
abierto, recurre a pedir \texttt{/api/live} por HTTP cada 1.5\,s. El
criterio es la robustez: un WebSocket puede caerse por muchas razones ajenas
al firmware (el navegador durmiendo la pestaña, un cambio de red en el
celular), y el panel no debería quedar congelado solo porque ese canal se
cortó.

\subsection{JSON con ArduinoJson}
Las respuestas HTTP y los mensajes WebSocket se serializan con la librería
ArduinoJson \cite{arduinojson2024} en vez de concatenar \textit{strings} a
mano. El criterio: construir JSON por concatenación de texto es una fuente
común de errores de escape/formato difíciles de depurar en un
microcontrolador, mientras que ArduinoJson gestiona el documento en un
búfer de tamaño acotado y conocido, algo relevante en un sistema con 320\,KB
de RAM total repartidos entre cinco tareas.

\begin{table}[H]
\centering
\caption{Endpoints servidos por ESPAsyncWebServer (\texttt{src/web\_server.cpp}).}
\begin{tabular}{@{}lll@{}}
\toprule
Método & Ruta & Descripción \\
\midrule
GET  & \texttt{/}              & Panel del dashboard (HTML embebido en el firmware) \\
WS   & \texttt{/ws}             & WebSocket, empuja el estado en vivo cada 500\,ms \\
GET  & \texttt{/api/live}       & Snapshot JSON del estado actual \\
GET  & \texttt{/api/files}      & Lista de archivos en \texttt{/logs} \\
GET  & \texttt{/api/download}   & Descarga un CSV completo \\
GET  & \texttt{/api/chartdata}  & Puntos decimados para graficar \\
POST & \texttt{/api/reset}      & Reinicia el contador de viaje \\
\bottomrule
\end{tabular}
\end{table}

\subsection{Interfaz}
El HTML/CSS/JS está embebido por completo en el firmware, sin dependencias
de CDN externo —criterio directo del contexto de uso: el SoftAP no tiene
salida a internet, así que cualquier librería externa de gráficas
simplemente no cargaría—. El panel tiene dos pestañas (\textit{En vivo} e
\textit{Historial}), tarjetas destacadas con el consumo instantáneo,
litros y costo del viaje, casillas para elegir qué variables graficar, y un
\textit{tooltip} interactivo con línea guía. Cuando hay 2 o más variables
seleccionadas a la vez, cada una se normaliza a su propio rango (0-100\,\%)
para poder comparar formas de curva con unidades muy distintas (RPM 0-7000
frente a L/100km 0-20) sin que una aplaste visualmente a la otra.

% ============================================================
\section{Compilación y programación}
\label{sec:build}
% ============================================================

\subsection{PlatformIO}

\begin{figure}[H]
\begin{lstlisting}[style=plain,caption={Configuración relevante de \texttt{platformio.ini}.}]
board = esp32-s3-devkitc-1
board_build.flash_mode = qio
board_build.arduino.memory_type = qio_opi
board_build.psram_type = opi
board_upload.flash_size = 16MB
\end{lstlisting}
\end{figure}

Para compilar y subir: \texttt{pio run -t upload} y luego
\texttt{pio device monitor -b 115200}.

\subsection{USB nativo frente a UART externo}
El \textit{devkit} de desarrollo solo expone el \textbf{USB nativo} del
ESP32-S3 (el periférico USB-Serial-JTAG integrado, sin un chip puente
CP2102/CH340 aparte). Esto tiene dos implicancias para una PCB propia:

\begin{enumerate}
  \item Programar por UART \textbf{siempre funciona}, sin importar la
        configuración del firmware: \texttt{esptool} habla directo con el
        \textit{bootloader} de ROM del chip, independiente de cualquier
        opción de compilación de la aplicación.
  \item El firmware define \texttt{-DARDUINO\_USB\_CDC\_ON\_BOOT=1}. Con ese
        \textit{flag}, el objeto \texttt{Serial} de Arduino se enruta al
        USB nativo en vez del UART0 físico (GPIO43/44). Si la PCB propia usa
        un adaptador UART externo y no tiene el USB nativo conectado a
        nada, ese \textit{flag} debe quitarse —de lo contrario el monitor
        serie del adaptador UART no mostrará nada, aunque el firmware esté
        corriendo correctamente.
\end{enumerate}

\subsection{Circuito mínimo para programar/arrancar}
Cualquier placa propia con el ESP32-S3 necesita, como mínimo: \texttt{EN}
(reset) con \textit{pull-up} a 3.3\,V y un capacitor de $\sim$1\,$\mu$F a
GND; \texttt{GPIO0} (BOOT) con \textit{pull-up} a 3.3\,V y un botón a GND
para forzar el modo \textit{bootloader} manualmente; y, opcionalmente, un
circuito de auto-reset (2 transistores desde DTR/RTS del adaptador UART
hacia EN y GPIO0) para que PlatformIO entre solo al modo de programación sin
botones.

% ============================================================
\section{Solución de problemas comunes}
% ============================================================

\begin{itemize}
  \item \textbf{La subida se queda en "Connecting..." indefinidamente.}
        Entrar manualmente en modo \textit{bootloader} (mantener BOOT,
        pulsar y soltar RESET/EN, soltar BOOT, y recién ahí lanzar la
        subida). Verificar también que ningún otro programa tenga el
        puerto COM abierto.
  \item \textbf{El monitor serie no muestra nada tras programar.} Revisar
        que \texttt{ARDUINO\_USB\_CDC\_ON\_BOOT} coincida con el hardware
        real (Sección~\ref{sec:build}).
  \item \textbf{\texttt{sdSelectCard(): Select Failed} / \texttt{SD=0} en el
        arranque.} Casi siempre es cableado/alimentación: revisar tarjeta
        insertada, tensión correcta del módulo lector, continuidad de los 4
        pines SPI + GND, formato FAT32, y el \textit{pull-up} de
        10\,k$\Omega$ en CS si se soldó directo sin módulo.
  \item \textbf{El panel web no carga en el celular.} Escribir
        \texttt{http://192.168.4.1} con el esquema explícito —varios
        navegadores intentan \texttt{https://} por defecto y el dispositivo
        no tiene TLS—. Algunos celulares además desconectan una red WiFi
        sin salida a internet a menos que se les indique quedarse
        conectados igual.
\end{itemize}
 (o copia/pega el
% contenido) dentro de tu documento original, en el punto donde va este
% capítulo. Ajusta el nivel de \section por \subsection si en tu
% documento esto ya cuelga de un \chapter{Diseño del software}.
%
% Requisitos que van en el PREAMBULO de tu documento maestro (no aquí):
%
%   \usepackage{amsmath,amssymb}
%   \usepackage{booktabs}
%   \usepackage{array}
%   \usepackage{xcolor}
%   \usepackage{listings}
%   \usepackage{float}
%   \usepackage{lmodern}    % fuentes escalables -- requisito de microtype
%                           % si tu documento usa babel en espanol (que
%                           % carga fuentes EC de mapa de bits sin esto)
%   \usepackage{microtype}  % reduce cajas "overfull" en parrafos con
%                           % identificadores \texttt largos (p.ej.
%                           % FUEL_CALIBRATION_FACTOR); combinar con:
%   \tolerance=1500
%   \emergencystretch=3em
%
%   \definecolor{codebg}{RGB}{245,245,245}
%   \definecolor{codekw}{RGB}{0,90,170}
%   \definecolor{codecom}{RGB}{110,110,110}
%   \definecolor{codestr}{RGB}{160,40,40}
%   \lstdefinestyle{cpp}{
%     language=C++, backgroundcolor=\color{codebg}, basicstyle=\ttfamily\small,
%     keywordstyle=\color{codekw}\bfseries, commentstyle=\color{codecom}\itshape,
%     stringstyle=\color{codestr}, numbers=left, numberstyle=\tiny\color{gray},
%     numbersep=8pt, breaklines=true, showstringspaces=false,
%     frame=single, rulecolor=\color{gray!40}, tabsize=2,
%   }
%   \lstdefinestyle{plain}{
%     backgroundcolor=\color{codebg}, basicstyle=\ttfamily\small,
%     breaklines=true, showstringspaces=false,
%     frame=single, rulecolor=\color{gray!40},
%   }
%   \lstset{style=cpp}
%
% Bibliografía: las fuentes citadas aquí (\cite{...}) están en
% referencias.bib, para fusionar con la bibliografía única de tu
% documento (donde ya tengas \bibliography{...} o \addbibresource{...}).
% \cite{} es el comando base de LaTeX, compatible con bibliografía clásica
% (\bibliographystyle{apalike} da formato "(Autor, Año)" sin paquetes
% extra), con natbib, y con biblatex.
% ============================================================

Esta sección retoma el capítulo de diseño de software y describe, parte por
parte, el firmware que corre en el ESP32-S3: las herramientas con las que se
construyó, cómo reparte el trabajo entre tareas, el protocolo con el que
habla con el vehículo, el método de cálculo del consumo, cómo persiste los
datos, cómo expone un panel de control, y el criterio detrás de cada
elección de componente y de librería.

% ============================================================
\section{Herramientas y entorno de desarrollo}
\label{sec:herramientas}
% ============================================================

\subsection{Visual Studio Code + PlatformIO}
El firmware se desarrolla en Visual Studio Code con la extensión PlatformIO
IDE, en vez del IDE oficial de Arduino. La diferencia no es cosmética:

\begin{itemize}
  \item \textbf{Toolchain embebido integrado.} PlatformIO gestiona la
        descarga y versión exacta del framework
        (\texttt{framework-arduinoespressif32}) \cite{arduinoesp32core2024},
        el compilador cruzado (\texttt{toolchain-xtensa-esp32s3}) y
        \texttt{esptool} para programar, todo declarado y fijado en
        \texttt{platformio.ini} \cite{platformio2024}. Esto hace que el
        \textit{build} sea reproducible entre máquinas —un requisito real
        una vez que el proyecto tiene múltiples archivos fuente y
        dependencias externas—, algo que el modelo de "sketch único" del
        IDE de Arduino maneja peor.
  \item \textbf{Autocompletado real de C++.} Con un servidor de lenguaje
        indexando todo el árbol de \texttt{src/}/\texttt{include/}, moverse
        entre los siete archivos fuente del proyecto y sus cabeceras
        compartidas (\texttt{config.h}, \texttt{types.h}) es inmediato.
  \item \textbf{Terminal integrado.} Compilar, subir y abrir el monitor
        serie desde el mismo terminal donde se edita el código evita el
        cambio de contexto entre editor y herramienta externa, y deja el
        registro de errores de compilación con enlaces directos a archivo y
        línea.
\end{itemize}

\subsection{Por qué FreeRTOS}
Un sistema que a la vez sondea un bus CAN sensible a tiempos, refresca una
pantalla, escribe en una tarjeta SD y sirve un panel web por WiFi es, por
naturaleza, un problema \textbf{concurrente}. Hay dos formas razonables de
resolverlo: un único \texttt{loop()} con máquinas de estado manuales y
sondeo no bloqueante en cada punto donde se podría esperar —viable en
proyectos pequeños, pero frágil a medida que crecen—, o un sistema operativo
de tiempo real con prioridades y primitivas de sincronización reales. Este
proyecto usa la segunda opción.

FreeRTOS no es, en rigor, una dependencia externa que se "agregó" —es el
planificador que el propio núcleo Arduino-ESP32 ya usa por debajo
\cite{freertos2024}; usarlo explícitamente es aprovechar la plataforma en
vez de pelear contra ella con un \texttt{loop()} monolítico encima.
Concretamente, esto permite:

\begin{itemize}
  \item Darle al sondeo CAN la \textbf{prioridad más alta} de todas las
        tareas, garantizando que una transacción I2C lenta hacia el OLED o
        una descarga de CSV por WiFi nunca retrase la lectura del vehículo.
  \item Proteger el estado compartido con mutex reales en vez de banderas
        improvisadas.
  \item Que agregar una funcionalidad nueva sea, casi siempre, agregar una
        tarea nueva con una prioridad clara, en vez de insertar más lógica
        condicional dentro de un bucle ya complejo —el modo de bajo consumo
        de la Sección~\ref{sec:idle} solo requirió tocar el bucle de la
        tarea CAN y agregar una bandera de estado que las demás tareas ya
        sabían consultar.
\end{itemize}

\begin{figure}[H]
\begin{lstlisting}[caption={Creación de tareas FreeRTOS en \texttt{src/main.cpp}.}, label={fig:tasks}]
statusLedsInit();
xTaskCreatePinnedToCore(statusLedsTask, "leds", STACK_LED_TASK, nullptr,
                         PRIO_LED_TASK, nullptr, CORE_LED_TASK);

bool oledOk = displayInit();
bool sdOk   = sdLoggerInit();
bool canOk  = canObd2Init();

webServerInit(); // reporta ErrorCode::WIFI_FAIL si el softAP falla

if (canOk) {
  xTaskCreatePinnedToCore(canObd2Task, "canObd2", STACK_CAN_TASK, nullptr,
                           PRIO_CAN_TASK, nullptr, CORE_CAN_TASK);
}
if (oledOk) {
  xTaskCreatePinnedToCore(displayTask, "display", STACK_DISPLAY_TASK, nullptr,
                           PRIO_DISPLAY_TASK, nullptr, CORE_DISPLAY_TASK);
}
if (sdOk) {
  xTaskCreatePinnedToCore(sdLoggerTask, "sdLogger", STACK_SD_TASK, nullptr,
                           PRIO_SD_TASK, nullptr, CORE_SD_TASK);
}
\end{lstlisting}
\end{figure}

% ============================================================
\section{Interfaz con el hardware}
% ============================================================

El software está organizado alrededor de las interfaces físicas del
ESP32-S3 \cite{espressif2024datasheet,espressif2024trm}. Cada periférico se
eligió por una razón concreta, no por ser "lo que venía en el kit":

\begin{table}[H]
\centering
\caption{Mapa de pines definido en \texttt{include/config.h}.}
\begin{tabular}{@{}lll@{}}
\toprule
Bus & Señal & GPIO ESP32-S3 \\
\midrule
TWAI (CAN)   & TX $\rightarrow$ SN65HVD230 CTX & GPIO4 \\
TWAI (CAN)   & RX $\leftarrow$ SN65HVD230 CRX  & GPIO5 \\
I2C          & SDA $\rightarrow$ OLED          & GPIO8 \\
I2C          & SCL $\rightarrow$ OLED          & GPIO9 \\
SPI (SD)     & SCK                              & GPIO12 \\
SPI (SD)     & MISO                             & GPIO13 \\
SPI (SD)     & MOSI                             & GPIO11 \\
SPI (SD)     & CS                               & GPIO10 \\
LED          & CAN activo                       & GPIO15 \\
LED          & Actividad SD                     & GPIO16 \\
LED          & Servidor activo                  & GPIO17 \\
LED          & Error                            & GPIO18 \\
\bottomrule
\end{tabular}
\end{table}

\textit{Nota.} Pines reservados —no usar—: GPIO 26-37 (bus interno hacia el
flash y la PSRAM octal del módulo N16R8) y GPIO 0/3/19/20/43-46 (pines de
\textit{strapping}, USB nativo y consola UART0 del chip).

\subsection{Controlador CAN: TWAI integrado, no un chip externo}
El ESP32-S3 trae un controlador CAN (TWAI) integrado en el propio chip
\cite{espressif2024trm}. Esa es la razón por la que el diseño no usa un
controlador CAN externo por SPI como el MCP2515, común en proyectos sobre
Arduino Uno/Nano que no tienen esta capacidad de fábrica: usar el
controlador nativo ahorra un dispositivo, una interrupción y un bus SPI que
ya está ocupado por la microSD, y evita la latencia adicional de tener que
serializar cada trama CAN a través de otro periférico intermediario.

\subsection{Transceptor CAN: SN65HVD230}
El TWAI del ESP32-S3 solo genera niveles lógicos TTL; convertirlos a las
señales diferenciales CAN\_H/CAN\_L del bus físico requiere un transceptor.
Se eligió el SN65HVD230 \cite{sn65hvd230_datasheet} en concreto porque opera
nativamente a 3.3\,V —el mismo nivel lógico del ESP32-S3—, evitando un
\textit{level-shifter} que sí haría falta con transceptores clásicos de 5\,V
como el MCP2551. CAN\_H y CAN\_L se conectan a los pines 6 y 14 del conector
OBD2 (J1962).

\subsection{Pantalla: SSD1306 OLED por I2C}
Se eligió un OLED monocromático de $128\times64$ px con controlador SSD1306
\cite{solomon2008ssd1306} sobre I2C en vez de, por ejemplo, una pantalla TFT
a color por SPI, por tres razones: consume solo 2 pines (I2C) dejando el SPI
libre para la SD; su consumo (10-20\,mA encendida, y se puede apagar por
software, Sección~\ref{sec:idle}) es bajo comparado con un TFT retroiluminado;
y la resolución alcanza de sobra para los pocos números grandes que hace
falta mostrar al conductor —no se necesita color ni gráficos complejos en el
tablero físico, eso ya lo cubre el panel web.

\subsection{Almacenamiento: microSD por SPI dedicado}
La tarjeta microSD usa su propio bus SPI, sin compartirlo con ningún otro
periférico. La alternativa —reutilizar un bus SPI ya en uso— se descartó
porque el registro en la SD ocurre una vez por segundo de forma predecible
mientras que una eventual escritura de mayor tamaño (por ejemplo, servir una
descarga de CSV completa) no debe competir por el bus con otra
responsabilidad no relacionada; un bus dedicado simplifica el análisis de
tiempos y evita tener que coordinar dos periféricos distintos sobre el mismo
SPI.

% ============================================================
\section{Arquitectura de tareas}
% ============================================================

\begin{table}[H]
\centering
\caption{Tareas FreeRTOS, definidas en \texttt{include/config.h} y creadas en \texttt{src/main.cpp}.}
\begin{tabular}{@{}llccc@{}}
\toprule
Tarea & Función & Núcleo & Prioridad & Pila \\
\midrule
CAN/OBD2          & \texttt{canObd2Task}       & 0 & 5 & 4096 \\
SD logger         & \texttt{sdLoggerTask}      & 1 & 4 & 6144 \\
WebSocket push     & \texttt{webSocketPushTask} & 0 & 3 & 4096 \\
OLED display      & \texttt{displayTask}       & 1 & 2 & 4096 \\
LEDs de estado    & \texttt{statusLedsTask}    & 1 & 1 & 2048 \\
\bottomrule
\end{tabular}
\end{table}

El núcleo 0 concentra lo sensible a tiempos (sondeo CAN de alta prioridad +
la pila WiFi/AsyncTCP, que en Arduino-ESP32 corre en ese mismo núcleo). El
núcleo 1 queda con las tareas de interfaz/almacenamiento, que toleran más
\textit{jitter} sin afectar la lectura del vehículo. Las prioridades no son
arbitrarias: reflejan directamente qué tan caro es que cada tarea llegue
tarde (perder un ciclo de sondeo CAN corrompe el cálculo de consumo; llegar
tarde a refrescar un LED es imperceptible).

\subsection{Sincronización entre tareas}
Dos mutex protegen el estado compartido: \texttt{g\_dataMutex} protege
\texttt{g\_vehicleData} (la tarea CAN arma cada ciclo en una copia local y
solo toca la estructura global, bajo mutex, en el paso final de fusión —
Figura~\ref{fig:merge}); \texttt{g\_sdMutex} protege el bus SPI de la
microSD, compartido entre el \textit{logger} periódico y los endpoints de
descarga/gráfica del servidor web. El criterio para usar mutex en vez de,
por ejemplo, deshabilitar interrupciones o copiar datos sin protección, es
que ambos recursos (la estructura de datos y el bus SPI) tienen escritores y
lectores en tareas de distinta prioridad corriendo en núcleos distintos —un
mutex es la primitiva correcta quando el acceso puede, en principio,
bloquear brevemente sin romper los tiempos críticos del sistema.

\begin{figure}[H]
\begin{lstlisting}[caption={Fusión bajo mutex al final de cada ciclo, \texttt{src/can\_obd2.cpp}.}, label={fig:merge}]
if (xSemaphoreTake(g_dataMutex, pdMS_TO_TICKS(100)) == pdTRUE) {
  VehicleData &g = g_vehicleData;
  g.rpm = scratch.rpm;
  g.speed_kmh = scratch.speed_kmh;
  // ... resto de los campos ...
  g.can_status = status;
  g.data_valid = criticalOk;
  g.last_update_ms = now;

  if (criticalOk && dt > 0.0f && dt < 5.0f) {
    fuelCalcUpdate(g, dt);
  }
  xSemaphoreGive(g_dataMutex);
}
\end{lstlisting}
\end{figure}

\subsection{Manejo de errores}
\texttt{systemReportError(ErrorCode)} centraliza el reporte de fallos: guarda
el último código en \texttt{g\_systemStatus.error\_code} (dispara el parpadeo
del LED de error) y lo imprime por \texttt{Serial} con el detalle completo.
El criterio es separar dos audiencias distintas: el LED es para el
conductor, que solo necesita saber \textit{que algo falló}; el log serie es
para depuración en banco, donde sí importa el detalle completo y el
historial.

\begin{table}[H]
\centering
\caption{Códigos de \texttt{ErrorCode} (\texttt{include/types.h}).}
\begin{tabular}{@{}cl@{}}
\toprule
Código & Significado \\
\midrule
1 & Fallo al iniciar el driver TWAI/CAN \\
2 & Fallo al iniciar la microSD \\
3 & Fallo al iniciar el OLED \\
4 & Fallo al iniciar el punto de acceso WiFi \\
\bottomrule
\end{tabular}
\end{table}

% ============================================================
\section{Protocolo CAN / OBD2}
% ============================================================

El controlador TWAI del ESP32-S3 se configura a \textbf{500\,kbit/s}, el
estándar de la mayoría de vehículos livianos con OBD2 sobre CAN
(ISO 15765-4, \cite{iso157654}). Todas las
peticiones usadas caben en un único \textit{frame} CAN de 8 bytes
(\textit{single frame}), así que no hace falta implementar el mecanismo
multi-\textit{frame} de control de flujo ISO-TP —una simplificación
deliberada: ninguno de los PIDs que este proyecto necesita supera ese
tamaño, y soportar multi-\textit{frame} sin necesitarlo solo agrega
superficie de error.

\subsection{PIDs consultados (Modo 01)}
El conjunto de PIDs se eligió por ser el mínimo necesario para el cálculo
speed-density (Sección~\ref{sec:speeddensity}) más los que enriquecen el
registro sin costo relevante de tiempo de bus, siguiendo el estándar SAE
J1979 \cite{sae_j1979_2014}:

\begin{table}[H]
\centering
\caption{PIDs Mode 01 consultados por \texttt{canObd2Task}.}
\begin{tabular}{@{}cll@{}}
\toprule
PID & Dato & Uso \\
\midrule
\texttt{0x0C} & RPM                          & speed-density (crítico) \\
\texttt{0x0B} & MAP (kPa)                    & speed-density (crítico) \\
\texttt{0x0F} & IAT (\textdegree C)          & speed-density (crítico) \\
\texttt{0x0D} & Velocidad (km/h)             & L/100km, distancia (crítico) \\
\texttt{0x05} & Temperatura refrigerante     & mostrado en pantalla/CSV \\
\texttt{0x04} & Carga calculada (\%)         & mostrado en pantalla/CSV \\
\texttt{0x06} & Fuel trim corto plazo (\%)   & corrige el AFR estimado \\
\texttt{0x07} & Fuel trim largo plazo (\%)   & corrige el AFR estimado \\
\bottomrule
\end{tabular}
\end{table}

\textit{Nota.} Los cuatro PIDs marcados "crítico" son indispensables para el
cálculo speed-density; si la ECU no responde alguno, el ciclo completo se
marca inválido y no se integra al viaje.

\subsection{Formato de la petición}

\begin{figure}[H]
\begin{lstlisting}[style=plain,caption={Trama de petición Mode 01, PID genérico.}]
ID = 0x7DF
Data[0] = 0x02   ; SF, 2 bytes de payload (modo + pid)
Data[1] = 0x01   ; Modo 01: datos actuales
Data[2] = <PID>
Data[3..7] = 0x00 (relleno)
\end{lstlisting}
\end{figure}

Las ECUs responden en el rango \texttt{0x7E8}-\texttt{0x7EF}; se acepta la
primera respuesta cuyo \texttt{Data[1] == 0x41} y \texttt{Data[2]} coincida
con el PID pedido.

\subsection{Recuperación de bus-off}
Si el controlador entra en estado \texttt{BUS\_OFF}, la recuperación se
sondea en cada vuelta del bucle sin bloquear con una espera fija —según el
driver TWAI de ESP-IDF \cite{espressif2024trm}, la recuperación solo llega a
\texttt{STOPPED} tras contar 128 ocurrencias de 11 bits recesivos
consecutivos en el bus, un tiempo que depende del tráfico real y no se puede
predecir de antemano; sondear en vez de esperar un tiempo fijo evita tanto
reintentar demasiado pronto (fallar en vano) como esperar de más
innecesariamente.

% ============================================================
\section{Método de cálculo: Speed-Density}
\label{sec:speeddensity}
% ============================================================

El método speed-density estima el flujo másico de aire admitido por el motor
a partir de la ley de gases ideales, sin medirlo con un sensor MAF —la
razón de fondo para elegir este método sobre uno basado en MAF real es que
la mayoría de los vehículos con OBD2 reportan MAP, no flujo de aire directo,
así que speed-density es el único método viable sin instalar un sensor
adicional—:

\begin{equation}
\dot{m}_{aire} \;[\text{g/s}] =
\frac{\dfrac{VE}{100} \cdot RPM \cdot MAP\,[\text{kPa}] \cdot V_d\,[\text{L}] \cdot 28.97}
     {2 \cdot 8.314 \cdot T_{IAT}\,[\text{K}] \cdot 60}
\end{equation}

donde $VE$ es la eficiencia volumétrica (\%) interpolada de una tabla
calibrable (Sección~\ref{sec:ve}); $V_d$ la cilindrada en litros; $28.97$ la
masa molar del aire (g/mol); $8.314$ la constante universal de los gases
(L$\cdot$kPa/(mol$\cdot$K)); el factor $2$ corresponde a que un motor de 4
tiempos completa una admisión cada 2 vueltas de cigüeñal; y $T_{IAT}$ la
temperatura de admisión en Kelvin, con guarda de rango
[$-40$,\,$100$]\,\textdegree C.

\begin{figure}[H]
\begin{lstlisting}[caption={Núcleo del cálculo speed-density, \texttt{src/fuel\_calc.cpp}.}, label={fig:fuelcalc}]
float ve = fuelCalcVeForRpm(vd.rpm);
float iat_K = vd.iat_c + 273.15f;
if (iat_K < 233.15f || iat_K > 373.15f) iat_K = 288.15f;

const float MM_AIR = 28.97f; // g/mol
const float R = 8.314f;      // L*kPa/(mol*K)

float maf = (ve / 100.0f) * (float)vd.rpm * vd.map_kpa *
            ENGINE_DISPLACEMENT_L * MM_AIR /
            (2.0f * R * iat_K * 60.0f);
\end{lstlisting}
\end{figure}

\subsection{De aire a combustible}

\begin{equation}
\dot{m}_{comb}\,[\text{g/s}] = \frac{\dot{m}_{aire}}{AFR} \cdot k_{cal}
\end{equation}

$AFR$ es la relación aire/combustible, por defecto el valor estequiométrico
de la gasolina ($AFR_{stoich} = 14.7$), opcionalmente corregido con los
\textit{fuel trims} que reporta la ECU:

\begin{equation}
AFR_{efectivo} = \frac{AFR_{stoich}}{1 + \frac{STFT\% + LTFT\%}{100}}
\end{equation}

acotado a $[8, 25]$. $k_{cal}$ es \texttt{FUEL\_CALIBRATION\_FACTOR}
(Sección~\ref{sec:calibracion}), $1.0$ por defecto. El criterio para incluir
la corrección por \textit{fuel trim} es que, sin ella, el cálculo asume que
el motor siempre corre exactamente en estequiométrico, cosa que en la
práctica no ocurre; usar los \textit{trims} que la propia ECU ya reporta es
una corrección "gratis" —no requiere sensores adicionales— frente a un error
sistemático conocido.

\subsection{De masa a volumen}

\begin{align}
\dot{V}_{comb}\,[\text{L/h}] &= \frac{\dot{m}_{comb} \cdot 3600}{\rho_{comb}} \\[4pt]
L/100km &=
  \begin{cases}
    \dfrac{\dot{V}_{comb}}{v}\cdot 100 & \text{si } v > 2\,\text{km/h} \\[6pt]
    0 & \text{vehículo detenido}
  \end{cases}
\end{align}

con $\rho_{comb} = 745\,\text{g/L}$ y $v$ la velocidad en km/h.

\subsection{Integración del viaje}

\begin{align}
L_{viaje} &\mathrel{+}= \dot{m}_{comb} \cdot dt \,/\, \rho_{comb} \\
km_{viaje} &\mathrel{+}= v \cdot dt / 3600 \\
Bs_{viaje} &= L_{viaje} \cdot P_{L}
\end{align}

con $dt$ el paso de tiempo real entre ciclos —no un valor fijo, precisamente
para no perder precisión si un ciclo de sondeo tarda más o menos de lo
típico— y $P_L = 6.96$ (\texttt{FUEL\_PRICE\_BS\_PER\_L}) el precio de
referencia del litro.

\subsection{Tabla de eficiencia volumétrica (VE)}
\label{sec:ve}

\begin{table}[H]
\centering
\caption{\texttt{kVeTable}, \texttt{src/fuel\_calc.cpp}. Interpolación lineal entre puntos.}
\begin{tabular}{@{}cc@{}}
\toprule
RPM & VE (\%) \\
\midrule
700  & 60.0 \\
1500 & 74.0 \\
2500 & 85.0 \\
4000 & 88.0 \\
5500 & 83.0 \\
7000 & 74.0 \\
\bottomrule
\end{tabular}
\end{table}

\textit{Nota.} Esta tabla es un punto de partida genérico, no un dato de
fábrica del vehículo. Sin calibrarla, el consumo calculado tendrá un error
sistemático.

\subsection{Calibración contra un tanque real}
\label{sec:calibracion}
No existe forma de que el ESP32 sepa, por sí solo, si el número que calcula
es correcto —no hay con qué comparar dentro del propio microcontrolador—.
El procedimiento recomendado: llenar el tanque y anotar el odómetro, manejar
con normalidad hasta la próxima recarga, comparar los litros reales cargados
contra \texttt{trip\_fuel\_L} acumulado en esa misma distancia, y ajustar

\begin{equation}
\texttt{FUEL\_CALIBRATION\_FACTOR} = \frac{L_{real}}{L_{calculado}}
\end{equation}

en \texttt{config.h}. El criterio de aplicar un único factor multiplicativo
sobre \texttt{fuel\_gs}, en vez de reajustar toda la tabla VE a mano, es que
corrige de forma consistente el instantáneo, el acumulado del viaje y el
costo con un solo número, y es reversible/auditable —a diferencia de
retocar seis puntos de la tabla VE por prueba y error.

% ============================================================
\section{Registro en microSD}
% ============================================================

Cada arranque crea un archivo nuevo \texttt{/logs/trip\_NNN.csv} (índice
autoincremental calculado escaneando los archivos existentes, en vez de un
contador separado que podría desincronizarse tras un corte de energía). Se
escribe una fila por segundo, solo después del primer ciclo CAN válido, y se
pausa por completo durante el modo de bajo consumo (Sección~\ref{sec:idle}).

\subsection{Por qué CSV}
Se eligió CSV en vez de un formato binario compacto o JSON por fila por tres
razones: es directamente legible/abrible en cualquier hoja de cálculo sin
escribir un parser; el tamaño de fila (unos 100 bytes) es intrascendente
frente a la capacidad de una microSD moderna incluso a una fila por segundo
durante horas; y es el formato que el propio panel web ya necesita poder
descargar y volver a interpretar (Sección~\ref{sec:panelweb}), así que un
mismo formato sirve para ambos usos sin conversión intermedia.

\begin{figure}[H]
\begin{lstlisting}[style=plain,caption={Fila de ejemplo del CSV (encabezado + una fila de datos).}]
uptime_ms,rpm,speed_kmh,map_kpa,iat_c,coolant_c,ve_pct,maf_gs,
fuel_gs,instant_Lh,instant_L100km,trip_km,trip_fuel_L,
trip_avg_L100km,can_status,trip_cost_bs

184532,2150,62.0,58.0,29.0,89.0,81.4,7.82,0.532,1.91,3.09,
12.400,0.384,3.10,1,2.67
\end{lstlisting}
\end{figure}

\begin{table}[H]
\centering
\caption{Columnas escritas por \texttt{sdLoggerTask} (\texttt{src/sd\_logger.cpp}).}
\begin{tabular}{@{}ll@{}}
\toprule
Columna & Descripción \\
\midrule
\texttt{uptime\_ms}        & Milisegundos desde el arranque \\
\texttt{rpm}                & Revoluciones por minuto \\
\texttt{speed\_kmh}        & Velocidad, km/h \\
\texttt{map\_kpa}          & Presión absoluta del múltiple, kPa \\
\texttt{iat\_c}             & Temperatura del aire de admisión, \textdegree C \\
\texttt{coolant\_c}        & Temperatura del refrigerante, \textdegree C \\
\texttt{ve\_pct}            & Eficiencia volumétrica usada, \% \\
\texttt{maf\_gs}            & Flujo másico de aire calculado, g/s \\
\texttt{fuel\_gs}           & Flujo másico de combustible calculado, g/s \\
\texttt{instant\_Lh}       & Consumo instantáneo, L/h \\
\texttt{instant\_L100km}  & Consumo instantáneo, L/100km \\
\texttt{trip\_km}          & Distancia acumulada del viaje, km \\
\texttt{trip\_fuel\_L}    & Combustible acumulado del viaje, L \\
\texttt{trip\_avg\_L100km} & Promedio del viaje, L/100km \\
\texttt{can\_status}       & Estado del bus CAN (enum \texttt{CanStatus}) \\
\texttt{trip\_cost\_bs}    & Costo acumulado del viaje, Bs \\
\bottomrule
\end{tabular}
\end{table}

% ============================================================
\section{Modo de bajo consumo ("sin contacto")}
\label{sec:idle}
% ============================================================

\subsection{Motivación}
El dispositivo se alimenta del pin 16 del conector OBD2, que entrega
12\,V \textbf{constantes} sin importar la posición de la llave. Si el
firmware siguiera sondeando el bus a ritmo completo, con el panel web y
todos los periféricos encendidos las 24 horas, se convertiría en una carga
parásita permanente sobre la batería del vehículo. La primera versión de
este modo se basaba en detectar RPM$=0$ y solo bajaba el ritmo de sondeo
—una mejora real pero modesta, porque el ESP32-S3 seguía activo con el
WiFi encendido—. La versión actual va más lejos: separa la alimentación en
dos rieles y usa \textbf{deep sleep} real del ESP32-S3 mientras no hay
contacto.

\subsection{Dos rieles de alimentación}
La PCB entrega dos salidas de 3.3\,V independientes desde el 12\,V
constante del OBD2:

\begin{itemize}
  \item \textbf{Riel 1} (siempre encendido): ESP32-S3 + SN65HVD230. Tiene
        que quedar siempre alimentado porque el transceptor es lo único
        capaz de "escuchar" el bus CAN para saber cuándo volvió el
        contacto.
  \item \textbf{Riel 2} (conmutado por software): LEDs + OLED, y por
        separado la microSD. Se corta por completo mientras no hay
        contacto.
\end{itemize}

\subsection{Circuito de \textit{power-gating}: MOSFET y pines}
\label{sec:powergate}
Cada riel 2 se controla con un MOSFET de \textbf{canal P, como interruptor
de alto lado} (corta el 3.3\,V que entra al periférico, no su GND). Se
prefiere sobre un N-MOSFET de bajo lado porque con este último el GND del
periférico puede quedar ligeramente flotante mientras conduce, lo que
genera problemas de integridad de señal en buses referenciados a GND como
I2C (OLED) y SPI (microSD).

\begin{table}[H]
\centering
\caption{Pines de \textit{power-gating}, \texttt{include/config.h}.}
\begin{tabular}{@{}lll@{}}
\toprule
Señal & GPIO ESP32-S3 & Controla \\
\midrule
\texttt{PIN\_PWR\_LED\_OLED} & GPIO6 & LEDs + OLED \\
\texttt{PIN\_PWR\_SD}        & GPIO7 & microSD \\
\bottomrule
\end{tabular}
\end{table}

Ambos GPIO son de dominio RTC en el ESP32-S3 (rango GPIO0-21), requisito
para que \texttt{gpio\_hold\_en()} pueda retener su nivel lógico durante el
\textit{deep sleep} \cite{espressif2024trm}. Topología de cada MOSFET
(misma para los dos rieles):

\begin{figure}[H]
\begin{lstlisting}[style=plain,caption={Interruptor de alto lado con P-MOSFET, un canal por riel.}]
Riel 2 (3.3V) ---+--------------------- Source (P-MOSFET)
                 |
              [R_pullup ~10k]
                 |
GPIO --[R_gate ~1k]------------------- Gate
                                          |
                                        Drain --- carga (OLED+LED, o microSD)
\end{lstlisting}
\end{figure}

El criterio detrás de cada componente: como el GPIO y el riel 2 comparten
el mismo nivel lógico (3.3\,V), no hace falta un transistor adicional para
manejar la compuerta —GPIO en alto da $V_{gs}=0$ (MOSFET apagado); GPIO en
bajo da $V_{gs}\approx-3.3\,\text{V}$, muy por debajo del umbral típico de
un P-MOSFET de nivel lógico ($-1$ a $-2\,\text{V}$ aprox., p.\,ej. un
AO3401A o equivalente)—. La resistencia de \textit{pull-up} en la
compuerta asegura que el MOSFET arranque apagado por defecto durante el
breve instante antes de que el firmware configure el GPIO como salida; la
resistencia en serie hacia el GPIO limita la corriente de carga de la
compuerta.

\subsection{Cadena de regulación}
\label{sec:regchain}
La alimentación completa, del 12\,V del OBD2 a los dos rieles de 3.3\,V,
pasa por tres reguladores en cascada:

\begin{enumerate}
  \item \textbf{TPS54302}: convertidor \textit{buck} (conmutado) que baja
        el 12\,V del pin 16 a 5\,V. Es la única etapa de conversión
        conmutada del sistema; las dos siguientes son lineales.
  \item \textbf{AMS1117} (5\,V $\rightarrow$ 3.3\,V): alimenta el
        \textbf{riel 1} (ESP32-S3 + SN65HVD230). Nunca se apaga.
  \item \textbf{AP2112} (5\,V $\rightarrow$ 3.3\,V): alimenta el
        \textbf{riel 2}; su salida de 3.3\,V entra a los dos MOSFET de la
        Sección~\ref{sec:powergate}, que gatean OLED+LED y microSD por
        separado. El AP2112 en sí mismo nunca se apaga —solo se corta lo
        que hay \textit{después} de los MOSFET—, pero al ser una LDO de
        muy baja corriente de reposo su aporte al consumo en espera es
        pequeño de todos modos (Sección~\ref{sec:consumo}).
\end{enumerate}

Esta cadena (\textit{buck} + 2 LDO en paralelo desde el mismo 5\,V) es una
elección de diseño válida y común, pero tiene una consecuencia directa para
el consumo "sin contacto" que conviene tener presente: a diferencia del
AP2112, el \textbf{AMS1117 es una LDO de propósito general, no de bajo
consumo} —su corriente de reposo (\textit{ground current}) es
comparativamente alta, típicamente 5-8\,mA casi sin importar la carga en
su salida, con variación notable entre fabricantes clon del mismo chip—.
Como el AMS1117 alimenta justo el riel que nunca se apaga, esa corriente de
reposo se sostiene las 24\,horas del día, se detalla en la
Sección~\ref{sec:consumo}.

\subsection{Implementación en firmware}
\texttt{canObd2CheckContact()} manda un único \textit{ping} (PID
\texttt{0x0C}, RPM, obligatorio en todo ECU OBD2) y reporta si algo
respondió, sin importar el valor. Se usa en dos puntos:

\begin{itemize}
  \item Al arrancar (\texttt{setup()}, tanto en un encendido normal como al
        despertar del \textit{deep sleep}): si no hay contacto, el chip
        vuelve a dormir antes de tocar el riel 2 o el WiFi.
  \item Ya despierto, dentro de \texttt{canObd2Task}: si el bus deja de
        responder durante \texttt{CONTACT\_LOST\_DEBOUNCE\_CYCLES} ciclos
        seguidos (3, para no dormirse por un solo \textit{frame} perdido en
        un bus ocupado), se considera que se fue el contacto.
\end{itemize}

\begin{figure}[H]
\begin{lstlisting}[caption={Entrada a deep sleep, \texttt{src/power\_sleep.cpp}.}, label={fig:deepsleep}]
void enterDeepSleepUntilContact() {
  sdLoggerShutdown();  // cierra el archivo antes de cortarle la energia
  WiFi.mode(WIFI_OFF);

  powerGateSet(false);
  gpio_hold_en((gpio_num_t)PIN_PWR_LED_OLED);
  gpio_hold_en((gpio_num_t)PIN_PWR_SD);
  gpio_deep_sleep_hold_en(); // sin esto el hold no sobrevive al deep sleep

  esp_sleep_enable_timer_wakeup(DEEP_SLEEP_WAKE_INTERVAL_US);
  esp_deep_sleep_start();  // no retorna
}
\end{lstlisting}
\end{figure}

Cerrar el archivo de la SD \textbf{antes} de cortar su MOSFET es
obligatorio, no una formalidad: cortar la alimentación con un archivo
abierto arriesga corromper el sistema de archivos completo, no solo perder
la última fila. El \textit{deep sleep} borra toda la RAM salvo el dominio
RTC, así que al despertar \texttt{setup()} vuelve a correr desde el
principio, exactamente como un arranque en frío.

\subsection{Consumo estimado y autonomía de la batería}
\label{sec:consumo}
Los siguientes números son estimaciones a partir de valores típicos de
hoja de datos, no mediciones sobre la placa real —sirven para tener una
idea de orden de magnitud, no un valor exacto—.

\begin{table}[H]
\centering
\small
\caption{Contribuyentes al consumo promedio "sin contacto", cadena real (TPS54302 + AMS1117 + AP2112).}
\begin{tabular}{@{}p{1.7cm}p{3.6cm}p{1.7cm}p{4.8cm}@{}}
\toprule
Nodo & Fuente & Aporte típico & Nota \\
\midrule
3.3\,V (riel 1) & ESP32-S3 en deep sleep & $\sim$0.02\,mA & 10-25\,$\mu$A \cite{espressif2024datasheet} \\
3.3\,V (riel 1) & SN65HVD230, siempre activo & $\sim$5-10\,mA & Modo normal, no standby \cite{sn65hvd230_datasheet} \\
3.3\,V (riel 1) & Pulsos de despertar (cada 20\,s) & $\sim$1.5-2\,mA & Promediado: $\sim$150\,mA por $\sim$0.2-0.3\,s \\
5\,V & \textbf{AMS1117}, corriente de reposo propia & \textbf{$\sim$5-8\,mA} & Casi independiente de la carga \\
5\,V & AP2112, corriente de reposo propia & $\sim$0.06\,mA & Sin carga (MOSFET de salida abiertos) \\
12\,V & TPS54302, corriente de reposo propia & $\sim$0.1-0.3\,mA & \\
\midrule
\multicolumn{2}{@{}l}{\textbf{Total aproximado, lado 12\,V}} & \textbf{$\sim$7-9\,mA} & Ver conversión abajo \\
\bottomrule
\end{tabular}
\end{table}

El resultado es contraintuitivo pero importante: \textbf{el transceptor CAN
y la LDO del riel 1, no el ESP32-S3, dominan el consumo mientras el
dispositivo "duerme"} —porque ese riel no se puede cortar sin perder la
capacidad de detectar cuándo vuelve el contacto (Sección~\ref{sec:powergate}).
El \textit{deep sleep} resuelve la parte del ESP32-S3 casi por completo,
pero no puede resolver la del transceptor ni la de su LDO con esta
arquitectura de dos rieles fija.

Convertir la corriente del lado 3.3\,V/5\,V al lado 12\,V: el AMS1117 y el
AP2112 son LDO (no reducen corriente, solo voltaje, disipando la diferencia
como calor), así que la corriente que el TPS54302 debe entregar en su
salida de 5\,V es aproximadamente la suma de lo que consumen el riel 1
completo ($\sim$8.5\,mA de ESP32+SN65+pulsos) más la corriente de reposo
propia del AMS1117 ($\sim$6.5\,mA punto medio) y del AP2112
($\sim$0.06\,mA): $\approx$15\,mA a 5\,V. El TPS54302, al ser un
\textit{buck} conmutado, sí convierte potencia de forma más eficiente
($\sim$80\,\% en carga liviana, estimado):

\begin{equation}
I_{12V} \approx \frac{5\,\text{V} \times 15\,\text{mA}}{12\,\text{V} \times 0.80} + 0.2\,\text{mA (I}_q\text{ TPS54302)} \approx 8\,\text{mA}
\end{equation}

Tomando esos $\sim$8\,mA y una batería típica de auto de 50\,Ah (de la cual,
por buena práctica, no conviene usar más del 50\,\% para no acortar su vida
útil, es decir $\sim$25\,000\,mAh utilizables):

\begin{equation}
\frac{25\,000\,\text{mAh}}{8\,\text{mA}} \approx 3\,125\,\text{horas} \approx 130\,\text{días}
\end{equation}

Es decir, \textbf{el dispositivo por sí solo} tardaría del orden de 4-5
meses en consumir la mitad de una batería sana —menos que si el riel 1
usara una LDO de baja corriente de reposo en vez del AMS1117, pero todavía
un margen amplio frente al uso normal del vehículo—. En la práctica esto
importa menos de lo que parece, porque el propio vehículo ya tiene su
propio consumo parásito de fábrica (alarma, memorias de radio/ECU, reloj,
etc.), típicamente $\sim$20-50\,mA con todo "apagado" —del mismo orden o
mayor que el aporte de este dispositivo—, y ese consumo de base es el que
realmente determina cuánto aguanta la batería estacionada, con o sin el
ESP32-S3 conectado.

\textbf{Decisión de diseño:} se evaluó reemplazar el AMS1117 del riel 1 por
una LDO de baja corriente de reposo (el propio AP2112, un MCP1700, o una
TPS7A02 con $I_q$ en el orden de nanoamperios bajarían el consumo "sin
contacto" varios mA más), pero se optó por mantener el AMS1117 y el AP2112
tal como estaban ya elegidos en la placa —es el contribuyente individual
más grande después del SN65HVD230, y queda documentado aquí como una
compensación consciente entre simplicidad/costo de la lista de materiales y
la autonomía teórica de la batería, no un descuido.

\subsection{Recomendación de mantenimiento}
Con esos números, la recomendación práctica no cambia mucho respecto al
cuidado normal de cualquier batería de auto: \textbf{arrancar/circular el
vehículo al menos cada 2-3 semanas}, o usar un cargador de mantenimiento
(\textit{battery tender}) si va a quedar estacionado más tiempo que eso.
Con la cadena de reguladores actual (AMS1117 incluido) el dispositivo
agrega una fracción del consumo parásito que el auto ya tiene de fábrica
—notable, pero no dominante—, así que este intervalo sigue siendo una
recomendación razonable sin necesitar ser más estricta.

\subsection{¿Afecta consultar tan seguido?}
Con la arquitectura actual hay dos frecuencias distintas que vale la pena
separar:

\begin{itemize}
  \item \textbf{Sondeo CAN mientras hay contacto} (un ciclo completo de 8
        PIDs cada $\sim$0.2-1.4\,s): el bus CAN a 500\,kbit/s tiene
        capacidad de sobra para esto —cada trama de petición/respuesta son
        unos 100 bits, así que incluso varias por segundo representan bien
        menos del 1\,\% del ancho de banda del bus—. Una ECU real está
        diseñada para atender consultas de diagnóstico bastante más
        frecuentes que esto sin ningún problema; no hay riesgo de saturar
        el bus ni de sobrecargar la ECU.
  \item \textbf{Escritura a la microSD, una fila por segundo}
        (\texttt{SD\_LOG\_INTERVAL\_MS}): las escrituras son secuenciales
        —se van agregando al final del archivo, no se reescribe una y otra
        vez el mismo bloque—, que es el mejor caso posible para el
        desgaste de la memoria flash. Una fila ronda los 100 bytes; a una
        por segundo son unos 360\,KB por hora de viaje, insignificante
        frente a la resistencia típica de una microSD moderna (miles de
        ciclos de escritura por bloque, con \textit{wear leveling}
        integrado). No es un factor de desgaste relevante en la práctica.
\end{itemize}

% ============================================================
\section{Panel web}
\label{sec:panelweb}
% ============================================================

\subsection{Punto de acceso WiFi: por qué SoftAP en vez de solo STA}
El dispositivo siempre levanta su propio punto de acceso (SSID
\texttt{ESP32\_OBD2}, IP típica \texttt{192.168.4.1}) en vez de depender
únicamente de unirse a una red WiFi conocida. El criterio es el contexto de
uso real: dentro de un vehículo casi nunca hay una red WiFi conocida en
rango, así que un diseño solo-STA dejaría el panel inaccesible la mayor
parte del tiempo. El SoftAP garantiza acceso local sin esa dependencia;
unirse además como cliente a una red conocida queda como algo opcional,
únicamente para sincronizar hora por NTP cuando hay una disponible.

\subsection{Por qué ESPAsyncWebServer en vez del servidor síncrono de Arduino}
Se eligió ESPAsyncWebServer \cite{espasyncwebserver2024} sobre el
\texttt{WebServer} síncrono incluido en el core de Arduino-ESP32 porque
atiende las conexiones de forma asíncrona, basada en eventos: una petición
lenta (por ejemplo, una descarga de CSV grande) no bloquea el hilo mientras
espera que el cliente reciba los datos. Dado que el servidor corre en el
mismo núcleo que la pila WiFi (Sección~\ref{sec:herramientas}), un modelo
síncrono habría significado que una descarga lenta retrasara también el
resto del tráfico de red, incluyendo el WebSocket que empuja los datos en
vivo.

\subsection{Empuje en vivo: WebSocket con respaldo por \textit{polling}}
El estado se empuja a los clientes por WebSocket cada 500\,ms
—Figura~\ref{fig:merge} ya mostró de dónde sale ese dato—, con un
mecanismo de respaldo: si el navegador detecta que el WebSocket no está
abierto, recurre a pedir \texttt{/api/live} por HTTP cada 1.5\,s. El
criterio es la robustez: un WebSocket puede caerse por muchas razones ajenas
al firmware (el navegador durmiendo la pestaña, un cambio de red en el
celular), y el panel no debería quedar congelado solo porque ese canal se
cortó.

\subsection{JSON con ArduinoJson}
Las respuestas HTTP y los mensajes WebSocket se serializan con la librería
ArduinoJson \cite{arduinojson2024} en vez de concatenar \textit{strings} a
mano. El criterio: construir JSON por concatenación de texto es una fuente
común de errores de escape/formato difíciles de depurar en un
microcontrolador, mientras que ArduinoJson gestiona el documento en un
búfer de tamaño acotado y conocido, algo relevante en un sistema con 320\,KB
de RAM total repartidos entre cinco tareas.

\begin{table}[H]
\centering
\caption{Endpoints servidos por ESPAsyncWebServer (\texttt{src/web\_server.cpp}).}
\begin{tabular}{@{}lll@{}}
\toprule
Método & Ruta & Descripción \\
\midrule
GET  & \texttt{/}              & Panel del dashboard (HTML embebido en el firmware) \\
WS   & \texttt{/ws}             & WebSocket, empuja el estado en vivo cada 500\,ms \\
GET  & \texttt{/api/live}       & Snapshot JSON del estado actual \\
GET  & \texttt{/api/files}      & Lista de archivos en \texttt{/logs} \\
GET  & \texttt{/api/download}   & Descarga un CSV completo \\
GET  & \texttt{/api/chartdata}  & Puntos decimados para graficar \\
POST & \texttt{/api/reset}      & Reinicia el contador de viaje \\
\bottomrule
\end{tabular}
\end{table}

\subsection{Interfaz}
El HTML/CSS/JS está embebido por completo en el firmware, sin dependencias
de CDN externo —criterio directo del contexto de uso: el SoftAP no tiene
salida a internet, así que cualquier librería externa de gráficas
simplemente no cargaría—. El panel tiene dos pestañas (\textit{En vivo} e
\textit{Historial}), tarjetas destacadas con el consumo instantáneo,
litros y costo del viaje, casillas para elegir qué variables graficar, y un
\textit{tooltip} interactivo con línea guía. Cuando hay 2 o más variables
seleccionadas a la vez, cada una se normaliza a su propio rango (0-100\,\%)
para poder comparar formas de curva con unidades muy distintas (RPM 0-7000
frente a L/100km 0-20) sin que una aplaste visualmente a la otra.

% ============================================================
\section{Compilación y programación}
\label{sec:build}
% ============================================================

\subsection{PlatformIO}

\begin{figure}[H]
\begin{lstlisting}[style=plain,caption={Configuración relevante de \texttt{platformio.ini}.}]
board = esp32-s3-devkitc-1
board_build.flash_mode = qio
board_build.arduino.memory_type = qio_opi
board_build.psram_type = opi
board_upload.flash_size = 16MB
\end{lstlisting}
\end{figure}

Para compilar y subir: \texttt{pio run -t upload} y luego
\texttt{pio device monitor -b 115200}.

\subsection{USB nativo frente a UART externo}
El \textit{devkit} de desarrollo solo expone el \textbf{USB nativo} del
ESP32-S3 (el periférico USB-Serial-JTAG integrado, sin un chip puente
CP2102/CH340 aparte). Esto tiene dos implicancias para una PCB propia:

\begin{enumerate}
  \item Programar por UART \textbf{siempre funciona}, sin importar la
        configuración del firmware: \texttt{esptool} habla directo con el
        \textit{bootloader} de ROM del chip, independiente de cualquier
        opción de compilación de la aplicación.
  \item El firmware define \texttt{-DARDUINO\_USB\_CDC\_ON\_BOOT=1}. Con ese
        \textit{flag}, el objeto \texttt{Serial} de Arduino se enruta al
        USB nativo en vez del UART0 físico (GPIO43/44). Si la PCB propia usa
        un adaptador UART externo y no tiene el USB nativo conectado a
        nada, ese \textit{flag} debe quitarse —de lo contrario el monitor
        serie del adaptador UART no mostrará nada, aunque el firmware esté
        corriendo correctamente.
\end{enumerate}

\subsection{Circuito mínimo para programar/arrancar}
Cualquier placa propia con el ESP32-S3 necesita, como mínimo: \texttt{EN}
(reset) con \textit{pull-up} a 3.3\,V y un capacitor de $\sim$1\,$\mu$F a
GND; \texttt{GPIO0} (BOOT) con \textit{pull-up} a 3.3\,V y un botón a GND
para forzar el modo \textit{bootloader} manualmente; y, opcionalmente, un
circuito de auto-reset (2 transistores desde DTR/RTS del adaptador UART
hacia EN y GPIO0) para que PlatformIO entre solo al modo de programación sin
botones.

% ============================================================
\section{Solución de problemas comunes}
% ============================================================

\begin{itemize}
  \item \textbf{La subida se queda en "Connecting..." indefinidamente.}
        Entrar manualmente en modo \textit{bootloader} (mantener BOOT,
        pulsar y soltar RESET/EN, soltar BOOT, y recién ahí lanzar la
        subida). Verificar también que ningún otro programa tenga el
        puerto COM abierto.
  \item \textbf{El monitor serie no muestra nada tras programar.} Revisar
        que \texttt{ARDUINO\_USB\_CDC\_ON\_BOOT} coincida con el hardware
        real (Sección~\ref{sec:build}).
  \item \textbf{\texttt{sdSelectCard(): Select Failed} / \texttt{SD=0} en el
        arranque.} Casi siempre es cableado/alimentación: revisar tarjeta
        insertada, tensión correcta del módulo lector, continuidad de los 4
        pines SPI + GND, formato FAT32, y el \textit{pull-up} de
        10\,k$\Omega$ en CS si se soldó directo sin módulo.
  \item \textbf{El panel web no carga en el celular.} Escribir
        \texttt{http://192.168.4.1} con el esquema explícito —varios
        navegadores intentan \texttt{https://} por defecto y el dispositivo
        no tiene TLS—. Algunos celulares además desconectan una red WiFi
        sin salida a internet a menos que se les indique quedarse
        conectados igual.
\end{itemize}
 (o copia/pega el
% contenido) dentro de tu documento original, en el punto donde va este
% capítulo. Ajusta el nivel de \section por \subsection si en tu
% documento esto ya cuelga de un \chapter{Diseño del software}.
%
% Requisitos que van en el PREAMBULO de tu documento maestro (no aquí):
%
%   \usepackage{amsmath,amssymb}
%   \usepackage{booktabs}
%   \usepackage{array}
%   \usepackage{xcolor}
%   \usepackage{listings}
%   \usepackage{float}
%   \usepackage{lmodern}    % fuentes escalables -- requisito de microtype
%                           % si tu documento usa babel en espanol (que
%                           % carga fuentes EC de mapa de bits sin esto)
%   \usepackage{microtype}  % reduce cajas "overfull" en parrafos con
%                           % identificadores \texttt largos (p.ej.
%                           % FUEL_CALIBRATION_FACTOR); combinar con:
%   \tolerance=1500
%   \emergencystretch=3em
%
%   \definecolor{codebg}{RGB}{245,245,245}
%   \definecolor{codekw}{RGB}{0,90,170}
%   \definecolor{codecom}{RGB}{110,110,110}
%   \definecolor{codestr}{RGB}{160,40,40}
%   \lstdefinestyle{cpp}{
%     language=C++, backgroundcolor=\color{codebg}, basicstyle=\ttfamily\small,
%     keywordstyle=\color{codekw}\bfseries, commentstyle=\color{codecom}\itshape,
%     stringstyle=\color{codestr}, numbers=left, numberstyle=\tiny\color{gray},
%     numbersep=8pt, breaklines=true, showstringspaces=false,
%     frame=single, rulecolor=\color{gray!40}, tabsize=2,
%   }
%   \lstdefinestyle{plain}{
%     backgroundcolor=\color{codebg}, basicstyle=\ttfamily\small,
%     breaklines=true, showstringspaces=false,
%     frame=single, rulecolor=\color{gray!40},
%   }
%   \lstset{style=cpp}
%
% Bibliografía: las fuentes citadas aquí (\cite{...}) están en
% referencias.bib, para fusionar con la bibliografía única de tu
% documento (donde ya tengas \bibliography{...} o \addbibresource{...}).
% \cite{} es el comando base de LaTeX, compatible con bibliografía clásica
% (\bibliographystyle{apalike} da formato "(Autor, Año)" sin paquetes
% extra), con natbib, y con biblatex.
% ============================================================

Esta sección retoma el capítulo de diseño de software y describe, parte por
parte, el firmware que corre en el ESP32-S3: las herramientas con las que se
construyó, cómo reparte el trabajo entre tareas, el protocolo con el que
habla con el vehículo, el método de cálculo del consumo, cómo persiste los
datos, cómo expone un panel de control, y el criterio detrás de cada
elección de componente y de librería.

% ============================================================
\section{Herramientas y entorno de desarrollo}
\label{sec:herramientas}
% ============================================================

\subsection{Visual Studio Code + PlatformIO}
El firmware se desarrolla en Visual Studio Code con la extensión PlatformIO
IDE, en vez del IDE oficial de Arduino. La diferencia no es cosmética:

\begin{itemize}
  \item \textbf{Toolchain embebido integrado.} PlatformIO gestiona la
        descarga y versión exacta del framework
        (\texttt{framework-arduinoespressif32}) \cite{arduinoesp32core2024},
        el compilador cruzado (\texttt{toolchain-xtensa-esp32s3}) y
        \texttt{esptool} para programar, todo declarado y fijado en
        \texttt{platformio.ini} \cite{platformio2024}. Esto hace que el
        \textit{build} sea reproducible entre máquinas —un requisito real
        una vez que el proyecto tiene múltiples archivos fuente y
        dependencias externas—, algo que el modelo de "sketch único" del
        IDE de Arduino maneja peor.
  \item \textbf{Autocompletado real de C++.} Con un servidor de lenguaje
        indexando todo el árbol de \texttt{src/}/\texttt{include/}, moverse
        entre los siete archivos fuente del proyecto y sus cabeceras
        compartidas (\texttt{config.h}, \texttt{types.h}) es inmediato.
  \item \textbf{Terminal integrado.} Compilar, subir y abrir el monitor
        serie desde el mismo terminal donde se edita el código evita el
        cambio de contexto entre editor y herramienta externa, y deja el
        registro de errores de compilación con enlaces directos a archivo y
        línea.
\end{itemize}

\subsection{Por qué FreeRTOS}
Un sistema que a la vez sondea un bus CAN sensible a tiempos, refresca una
pantalla, escribe en una tarjeta SD y sirve un panel web por WiFi es, por
naturaleza, un problema \textbf{concurrente}. Hay dos formas razonables de
resolverlo: un único \texttt{loop()} con máquinas de estado manuales y
sondeo no bloqueante en cada punto donde se podría esperar —viable en
proyectos pequeños, pero frágil a medida que crecen—, o un sistema operativo
de tiempo real con prioridades y primitivas de sincronización reales. Este
proyecto usa la segunda opción.

FreeRTOS no es, en rigor, una dependencia externa que se "agregó" —es el
planificador que el propio núcleo Arduino-ESP32 ya usa por debajo
\cite{freertos2024}; usarlo explícitamente es aprovechar la plataforma en
vez de pelear contra ella con un \texttt{loop()} monolítico encima.
Concretamente, esto permite:

\begin{itemize}
  \item Darle al sondeo CAN la \textbf{prioridad más alta} de todas las
        tareas, garantizando que una transacción I2C lenta hacia el OLED o
        una descarga de CSV por WiFi nunca retrase la lectura del vehículo.
  \item Proteger el estado compartido con mutex reales en vez de banderas
        improvisadas.
  \item Que agregar una funcionalidad nueva sea, casi siempre, agregar una
        tarea nueva con una prioridad clara, en vez de insertar más lógica
        condicional dentro de un bucle ya complejo —el modo de bajo consumo
        de la Sección~\ref{sec:idle} solo requirió tocar el bucle de la
        tarea CAN y agregar una bandera de estado que las demás tareas ya
        sabían consultar.
\end{itemize}

\begin{figure}[H]
\begin{lstlisting}[caption={Creación de tareas FreeRTOS en \texttt{src/main.cpp}.}, label={fig:tasks}]
statusLedsInit();
xTaskCreatePinnedToCore(statusLedsTask, "leds", STACK_LED_TASK, nullptr,
                         PRIO_LED_TASK, nullptr, CORE_LED_TASK);

bool oledOk = displayInit();
bool sdOk   = sdLoggerInit();
bool canOk  = canObd2Init();

webServerInit(); // reporta ErrorCode::WIFI_FAIL si el softAP falla

if (canOk) {
  xTaskCreatePinnedToCore(canObd2Task, "canObd2", STACK_CAN_TASK, nullptr,
                           PRIO_CAN_TASK, nullptr, CORE_CAN_TASK);
}
if (oledOk) {
  xTaskCreatePinnedToCore(displayTask, "display", STACK_DISPLAY_TASK, nullptr,
                           PRIO_DISPLAY_TASK, nullptr, CORE_DISPLAY_TASK);
}
if (sdOk) {
  xTaskCreatePinnedToCore(sdLoggerTask, "sdLogger", STACK_SD_TASK, nullptr,
                           PRIO_SD_TASK, nullptr, CORE_SD_TASK);
}
\end{lstlisting}
\end{figure}

% ============================================================
\section{Interfaz con el hardware}
% ============================================================

El software está organizado alrededor de las interfaces físicas del
ESP32-S3 \cite{espressif2024datasheet,espressif2024trm}. Cada periférico se
eligió por una razón concreta, no por ser "lo que venía en el kit":

\begin{table}[H]
\centering
\caption{Mapa de pines definido en \texttt{include/config.h}.}
\begin{tabular}{@{}lll@{}}
\toprule
Bus & Señal & GPIO ESP32-S3 \\
\midrule
TWAI (CAN)   & TX $\rightarrow$ SN65HVD230 CTX & GPIO4 \\
TWAI (CAN)   & RX $\leftarrow$ SN65HVD230 CRX  & GPIO5 \\
I2C          & SDA $\rightarrow$ OLED          & GPIO8 \\
I2C          & SCL $\rightarrow$ OLED          & GPIO9 \\
SPI (SD)     & SCK                              & GPIO12 \\
SPI (SD)     & MISO                             & GPIO13 \\
SPI (SD)     & MOSI                             & GPIO11 \\
SPI (SD)     & CS                               & GPIO10 \\
LED          & CAN activo                       & GPIO15 \\
LED          & Actividad SD                     & GPIO16 \\
LED          & Servidor activo                  & GPIO17 \\
LED          & Error                            & GPIO18 \\
\bottomrule
\end{tabular}
\end{table}

\textit{Nota.} Pines reservados —no usar—: GPIO 26-37 (bus interno hacia el
flash y la PSRAM octal del módulo N16R8) y GPIO 0/3/19/20/43-46 (pines de
\textit{strapping}, USB nativo y consola UART0 del chip).

\subsection{Controlador CAN: TWAI integrado, no un chip externo}
El ESP32-S3 trae un controlador CAN (TWAI) integrado en el propio chip
\cite{espressif2024trm}. Esa es la razón por la que el diseño no usa un
controlador CAN externo por SPI como el MCP2515, común en proyectos sobre
Arduino Uno/Nano que no tienen esta capacidad de fábrica: usar el
controlador nativo ahorra un dispositivo, una interrupción y un bus SPI que
ya está ocupado por la microSD, y evita la latencia adicional de tener que
serializar cada trama CAN a través de otro periférico intermediario.

\subsection{Transceptor CAN: SN65HVD230}
El TWAI del ESP32-S3 solo genera niveles lógicos TTL; convertirlos a las
señales diferenciales CAN\_H/CAN\_L del bus físico requiere un transceptor.
Se eligió el SN65HVD230 \cite{sn65hvd230_datasheet} en concreto porque opera
nativamente a 3.3\,V —el mismo nivel lógico del ESP32-S3—, evitando un
\textit{level-shifter} que sí haría falta con transceptores clásicos de 5\,V
como el MCP2551. CAN\_H y CAN\_L se conectan a los pines 6 y 14 del conector
OBD2 (J1962).

\subsection{Pantalla: SSD1306 OLED por I2C}
Se eligió un OLED monocromático de $128\times64$ px con controlador SSD1306
\cite{solomon2008ssd1306} sobre I2C en vez de, por ejemplo, una pantalla TFT
a color por SPI, por tres razones: consume solo 2 pines (I2C) dejando el SPI
libre para la SD; su consumo (10-20\,mA encendida, y se puede apagar por
software, Sección~\ref{sec:idle}) es bajo comparado con un TFT retroiluminado;
y la resolución alcanza de sobra para los pocos números grandes que hace
falta mostrar al conductor —no se necesita color ni gráficos complejos en el
tablero físico, eso ya lo cubre el panel web.

\subsection{Almacenamiento: microSD por SPI dedicado}
La tarjeta microSD usa su propio bus SPI, sin compartirlo con ningún otro
periférico. La alternativa —reutilizar un bus SPI ya en uso— se descartó
porque el registro en la SD ocurre una vez por segundo de forma predecible
mientras que una eventual escritura de mayor tamaño (por ejemplo, servir una
descarga de CSV completa) no debe competir por el bus con otra
responsabilidad no relacionada; un bus dedicado simplifica el análisis de
tiempos y evita tener que coordinar dos periféricos distintos sobre el mismo
SPI.

% ============================================================
\section{Arquitectura de tareas}
% ============================================================

\begin{table}[H]
\centering
\caption{Tareas FreeRTOS, definidas en \texttt{include/config.h} y creadas en \texttt{src/main.cpp}.}
\begin{tabular}{@{}llccc@{}}
\toprule
Tarea & Función & Núcleo & Prioridad & Pila \\
\midrule
CAN/OBD2          & \texttt{canObd2Task}       & 0 & 5 & 4096 \\
SD logger         & \texttt{sdLoggerTask}      & 1 & 4 & 6144 \\
WebSocket push     & \texttt{webSocketPushTask} & 0 & 3 & 4096 \\
OLED display      & \texttt{displayTask}       & 1 & 2 & 4096 \\
LEDs de estado    & \texttt{statusLedsTask}    & 1 & 1 & 2048 \\
\bottomrule
\end{tabular}
\end{table}

El núcleo 0 concentra lo sensible a tiempos (sondeo CAN de alta prioridad +
la pila WiFi/AsyncTCP, que en Arduino-ESP32 corre en ese mismo núcleo). El
núcleo 1 queda con las tareas de interfaz/almacenamiento, que toleran más
\textit{jitter} sin afectar la lectura del vehículo. Las prioridades no son
arbitrarias: reflejan directamente qué tan caro es que cada tarea llegue
tarde (perder un ciclo de sondeo CAN corrompe el cálculo de consumo; llegar
tarde a refrescar un LED es imperceptible).

\subsection{Sincronización entre tareas}
Dos mutex protegen el estado compartido: \texttt{g\_dataMutex} protege
\texttt{g\_vehicleData} (la tarea CAN arma cada ciclo en una copia local y
solo toca la estructura global, bajo mutex, en el paso final de fusión —
Figura~\ref{fig:merge}); \texttt{g\_sdMutex} protege el bus SPI de la
microSD, compartido entre el \textit{logger} periódico y los endpoints de
descarga/gráfica del servidor web. El criterio para usar mutex en vez de,
por ejemplo, deshabilitar interrupciones o copiar datos sin protección, es
que ambos recursos (la estructura de datos y el bus SPI) tienen escritores y
lectores en tareas de distinta prioridad corriendo en núcleos distintos —un
mutex es la primitiva correcta quando el acceso puede, en principio,
bloquear brevemente sin romper los tiempos críticos del sistema.

\begin{figure}[H]
\begin{lstlisting}[caption={Fusión bajo mutex al final de cada ciclo, \texttt{src/can\_obd2.cpp}.}, label={fig:merge}]
if (xSemaphoreTake(g_dataMutex, pdMS_TO_TICKS(100)) == pdTRUE) {
  VehicleData &g = g_vehicleData;
  g.rpm = scratch.rpm;
  g.speed_kmh = scratch.speed_kmh;
  // ... resto de los campos ...
  g.can_status = status;
  g.data_valid = criticalOk;
  g.last_update_ms = now;

  if (criticalOk && dt > 0.0f && dt < 5.0f) {
    fuelCalcUpdate(g, dt);
  }
  xSemaphoreGive(g_dataMutex);
}
\end{lstlisting}
\end{figure}

\subsection{Manejo de errores}
\texttt{systemReportError(ErrorCode)} centraliza el reporte de fallos: guarda
el último código en \texttt{g\_systemStatus.error\_code} (dispara el parpadeo
del LED de error) y lo imprime por \texttt{Serial} con el detalle completo.
El criterio es separar dos audiencias distintas: el LED es para el
conductor, que solo necesita saber \textit{que algo falló}; el log serie es
para depuración en banco, donde sí importa el detalle completo y el
historial.

\begin{table}[H]
\centering
\caption{Códigos de \texttt{ErrorCode} (\texttt{include/types.h}).}
\begin{tabular}{@{}cl@{}}
\toprule
Código & Significado \\
\midrule
1 & Fallo al iniciar el driver TWAI/CAN \\
2 & Fallo al iniciar la microSD \\
3 & Fallo al iniciar el OLED \\
4 & Fallo al iniciar el punto de acceso WiFi \\
\bottomrule
\end{tabular}
\end{table}

% ============================================================
\section{Protocolo CAN / OBD2}
% ============================================================

El controlador TWAI del ESP32-S3 se configura a \textbf{500\,kbit/s}, el
estándar de la mayoría de vehículos livianos con OBD2 sobre CAN
(ISO 15765-4, \cite{iso157654}). Todas las
peticiones usadas caben en un único \textit{frame} CAN de 8 bytes
(\textit{single frame}), así que no hace falta implementar el mecanismo
multi-\textit{frame} de control de flujo ISO-TP —una simplificación
deliberada: ninguno de los PIDs que este proyecto necesita supera ese
tamaño, y soportar multi-\textit{frame} sin necesitarlo solo agrega
superficie de error.

\subsection{PIDs consultados (Modo 01)}
El conjunto de PIDs se eligió por ser el mínimo necesario para el cálculo
speed-density (Sección~\ref{sec:speeddensity}) más los que enriquecen el
registro sin costo relevante de tiempo de bus, siguiendo el estándar SAE
J1979 \cite{sae_j1979_2014}:

\begin{table}[H]
\centering
\caption{PIDs Mode 01 consultados por \texttt{canObd2Task}.}
\begin{tabular}{@{}cll@{}}
\toprule
PID & Dato & Uso \\
\midrule
\texttt{0x0C} & RPM                          & speed-density (crítico) \\
\texttt{0x0B} & MAP (kPa)                    & speed-density (crítico) \\
\texttt{0x0F} & IAT (\textdegree C)          & speed-density (crítico) \\
\texttt{0x0D} & Velocidad (km/h)             & L/100km, distancia (crítico) \\
\texttt{0x05} & Temperatura refrigerante     & mostrado en pantalla/CSV \\
\texttt{0x04} & Carga calculada (\%)         & mostrado en pantalla/CSV \\
\texttt{0x06} & Fuel trim corto plazo (\%)   & corrige el AFR estimado \\
\texttt{0x07} & Fuel trim largo plazo (\%)   & corrige el AFR estimado \\
\bottomrule
\end{tabular}
\end{table}

\textit{Nota.} Los cuatro PIDs marcados "crítico" son indispensables para el
cálculo speed-density; si la ECU no responde alguno, el ciclo completo se
marca inválido y no se integra al viaje.

\subsection{Formato de la petición}

\begin{figure}[H]
\begin{lstlisting}[style=plain,caption={Trama de petición Mode 01, PID genérico.}]
ID = 0x7DF
Data[0] = 0x02   ; SF, 2 bytes de payload (modo + pid)
Data[1] = 0x01   ; Modo 01: datos actuales
Data[2] = <PID>
Data[3..7] = 0x00 (relleno)
\end{lstlisting}
\end{figure}

Las ECUs responden en el rango \texttt{0x7E8}-\texttt{0x7EF}; se acepta la
primera respuesta cuyo \texttt{Data[1] == 0x41} y \texttt{Data[2]} coincida
con el PID pedido.

\subsection{Recuperación de bus-off}
Si el controlador entra en estado \texttt{BUS\_OFF}, la recuperación se
sondea en cada vuelta del bucle sin bloquear con una espera fija —según el
driver TWAI de ESP-IDF \cite{espressif2024trm}, la recuperación solo llega a
\texttt{STOPPED} tras contar 128 ocurrencias de 11 bits recesivos
consecutivos en el bus, un tiempo que depende del tráfico real y no se puede
predecir de antemano; sondear en vez de esperar un tiempo fijo evita tanto
reintentar demasiado pronto (fallar en vano) como esperar de más
innecesariamente.

% ============================================================
\section{Método de cálculo: Speed-Density}
\label{sec:speeddensity}
% ============================================================

El método speed-density estima el flujo másico de aire admitido por el motor
a partir de la ley de gases ideales, sin medirlo con un sensor MAF —la
razón de fondo para elegir este método sobre uno basado en MAF real es que
la mayoría de los vehículos con OBD2 reportan MAP, no flujo de aire directo,
así que speed-density es el único método viable sin instalar un sensor
adicional—:

\begin{equation}
\dot{m}_{aire} \;[\text{g/s}] =
\frac{\dfrac{VE}{100} \cdot RPM \cdot MAP\,[\text{kPa}] \cdot V_d\,[\text{L}] \cdot 28.97}
     {2 \cdot 8.314 \cdot T_{IAT}\,[\text{K}] \cdot 60}
\end{equation}

donde $VE$ es la eficiencia volumétrica (\%) interpolada de una tabla
calibrable (Sección~\ref{sec:ve}); $V_d$ la cilindrada en litros; $28.97$ la
masa molar del aire (g/mol); $8.314$ la constante universal de los gases
(L$\cdot$kPa/(mol$\cdot$K)); el factor $2$ corresponde a que un motor de 4
tiempos completa una admisión cada 2 vueltas de cigüeñal; y $T_{IAT}$ la
temperatura de admisión en Kelvin, con guarda de rango
[$-40$,\,$100$]\,\textdegree C.

\begin{figure}[H]
\begin{lstlisting}[caption={Núcleo del cálculo speed-density, \texttt{src/fuel\_calc.cpp}.}, label={fig:fuelcalc}]
float ve = fuelCalcVeForRpm(vd.rpm);
float iat_K = vd.iat_c + 273.15f;
if (iat_K < 233.15f || iat_K > 373.15f) iat_K = 288.15f;

const float MM_AIR = 28.97f; // g/mol
const float R = 8.314f;      // L*kPa/(mol*K)

float maf = (ve / 100.0f) * (float)vd.rpm * vd.map_kpa *
            ENGINE_DISPLACEMENT_L * MM_AIR /
            (2.0f * R * iat_K * 60.0f);
\end{lstlisting}
\end{figure}

\subsection{De aire a combustible}

\begin{equation}
\dot{m}_{comb}\,[\text{g/s}] = \frac{\dot{m}_{aire}}{AFR} \cdot k_{cal}
\end{equation}

$AFR$ es la relación aire/combustible, por defecto el valor estequiométrico
de la gasolina ($AFR_{stoich} = 14.7$), opcionalmente corregido con los
\textit{fuel trims} que reporta la ECU:

\begin{equation}
AFR_{efectivo} = \frac{AFR_{stoich}}{1 + \frac{STFT\% + LTFT\%}{100}}
\end{equation}

acotado a $[8, 25]$. $k_{cal}$ es \texttt{FUEL\_CALIBRATION\_FACTOR}
(Sección~\ref{sec:calibracion}), $1.0$ por defecto. El criterio para incluir
la corrección por \textit{fuel trim} es que, sin ella, el cálculo asume que
el motor siempre corre exactamente en estequiométrico, cosa que en la
práctica no ocurre; usar los \textit{trims} que la propia ECU ya reporta es
una corrección "gratis" —no requiere sensores adicionales— frente a un error
sistemático conocido.

\subsection{De masa a volumen}

\begin{align}
\dot{V}_{comb}\,[\text{L/h}] &= \frac{\dot{m}_{comb} \cdot 3600}{\rho_{comb}} \\[4pt]
L/100km &=
  \begin{cases}
    \dfrac{\dot{V}_{comb}}{v}\cdot 100 & \text{si } v > 2\,\text{km/h} \\[6pt]
    0 & \text{vehículo detenido}
  \end{cases}
\end{align}

con $\rho_{comb} = 745\,\text{g/L}$ y $v$ la velocidad en km/h.

\subsection{Integración del viaje}

\begin{align}
L_{viaje} &\mathrel{+}= \dot{m}_{comb} \cdot dt \,/\, \rho_{comb} \\
km_{viaje} &\mathrel{+}= v \cdot dt / 3600 \\
Bs_{viaje} &= L_{viaje} \cdot P_{L}
\end{align}

con $dt$ el paso de tiempo real entre ciclos —no un valor fijo, precisamente
para no perder precisión si un ciclo de sondeo tarda más o menos de lo
típico— y $P_L = 6.96$ (\texttt{FUEL\_PRICE\_BS\_PER\_L}) el precio de
referencia del litro.

\subsection{Tabla de eficiencia volumétrica (VE)}
\label{sec:ve}

\begin{table}[H]
\centering
\caption{\texttt{kVeTable}, \texttt{src/fuel\_calc.cpp}. Interpolación lineal entre puntos.}
\begin{tabular}{@{}cc@{}}
\toprule
RPM & VE (\%) \\
\midrule
700  & 60.0 \\
1500 & 74.0 \\
2500 & 85.0 \\
4000 & 88.0 \\
5500 & 83.0 \\
7000 & 74.0 \\
\bottomrule
\end{tabular}
\end{table}

\textit{Nota.} Esta tabla es un punto de partida genérico, no un dato de
fábrica del vehículo. Sin calibrarla, el consumo calculado tendrá un error
sistemático.

\subsection{Calibración contra un tanque real}
\label{sec:calibracion}
No existe forma de que el ESP32 sepa, por sí solo, si el número que calcula
es correcto —no hay con qué comparar dentro del propio microcontrolador—.
El procedimiento recomendado: llenar el tanque y anotar el odómetro, manejar
con normalidad hasta la próxima recarga, comparar los litros reales cargados
contra \texttt{trip\_fuel\_L} acumulado en esa misma distancia, y ajustar

\begin{equation}
\texttt{FUEL\_CALIBRATION\_FACTOR} = \frac{L_{real}}{L_{calculado}}
\end{equation}

en \texttt{config.h}. El criterio de aplicar un único factor multiplicativo
sobre \texttt{fuel\_gs}, en vez de reajustar toda la tabla VE a mano, es que
corrige de forma consistente el instantáneo, el acumulado del viaje y el
costo con un solo número, y es reversible/auditable —a diferencia de
retocar seis puntos de la tabla VE por prueba y error.

% ============================================================
\section{Registro en microSD}
% ============================================================

Cada arranque crea un archivo nuevo \texttt{/logs/trip\_NNN.csv} (índice
autoincremental calculado escaneando los archivos existentes, en vez de un
contador separado que podría desincronizarse tras un corte de energía). Se
escribe una fila por segundo, solo después del primer ciclo CAN válido, y se
pausa por completo durante el modo de bajo consumo (Sección~\ref{sec:idle}).

\subsection{Por qué CSV}
Se eligió CSV en vez de un formato binario compacto o JSON por fila por tres
razones: es directamente legible/abrible en cualquier hoja de cálculo sin
escribir un parser; el tamaño de fila (unos 100 bytes) es intrascendente
frente a la capacidad de una microSD moderna incluso a una fila por segundo
durante horas; y es el formato que el propio panel web ya necesita poder
descargar y volver a interpretar (Sección~\ref{sec:panelweb}), así que un
mismo formato sirve para ambos usos sin conversión intermedia.

\begin{figure}[H]
\begin{lstlisting}[style=plain,caption={Fila de ejemplo del CSV (encabezado + una fila de datos).}]
uptime_ms,rpm,speed_kmh,map_kpa,iat_c,coolant_c,ve_pct,maf_gs,
fuel_gs,instant_Lh,instant_L100km,trip_km,trip_fuel_L,
trip_avg_L100km,can_status,trip_cost_bs

184532,2150,62.0,58.0,29.0,89.0,81.4,7.82,0.532,1.91,3.09,
12.400,0.384,3.10,1,2.67
\end{lstlisting}
\end{figure}

\begin{table}[H]
\centering
\caption{Columnas escritas por \texttt{sdLoggerTask} (\texttt{src/sd\_logger.cpp}).}
\begin{tabular}{@{}ll@{}}
\toprule
Columna & Descripción \\
\midrule
\texttt{uptime\_ms}        & Milisegundos desde el arranque \\
\texttt{rpm}                & Revoluciones por minuto \\
\texttt{speed\_kmh}        & Velocidad, km/h \\
\texttt{map\_kpa}          & Presión absoluta del múltiple, kPa \\
\texttt{iat\_c}             & Temperatura del aire de admisión, \textdegree C \\
\texttt{coolant\_c}        & Temperatura del refrigerante, \textdegree C \\
\texttt{ve\_pct}            & Eficiencia volumétrica usada, \% \\
\texttt{maf\_gs}            & Flujo másico de aire calculado, g/s \\
\texttt{fuel\_gs}           & Flujo másico de combustible calculado, g/s \\
\texttt{instant\_Lh}       & Consumo instantáneo, L/h \\
\texttt{instant\_L100km}  & Consumo instantáneo, L/100km \\
\texttt{trip\_km}          & Distancia acumulada del viaje, km \\
\texttt{trip\_fuel\_L}    & Combustible acumulado del viaje, L \\
\texttt{trip\_avg\_L100km} & Promedio del viaje, L/100km \\
\texttt{can\_status}       & Estado del bus CAN (enum \texttt{CanStatus}) \\
\texttt{trip\_cost\_bs}    & Costo acumulado del viaje, Bs \\
\bottomrule
\end{tabular}
\end{table}

% ============================================================
\section{Modo de bajo consumo ("sin contacto")}
\label{sec:idle}
% ============================================================

\subsection{Motivación}
El dispositivo se alimenta del pin 16 del conector OBD2, que entrega
12\,V \textbf{constantes} sin importar la posición de la llave. Si el
firmware siguiera sondeando el bus a ritmo completo, con el panel web y
todos los periféricos encendidos las 24 horas, se convertiría en una carga
parásita permanente sobre la batería del vehículo. La primera versión de
este modo se basaba en detectar RPM$=0$ y solo bajaba el ritmo de sondeo
—una mejora real pero modesta, porque el ESP32-S3 seguía activo con el
WiFi encendido—. La versión actual va más lejos: separa la alimentación en
dos rieles y usa \textbf{deep sleep} real del ESP32-S3 mientras no hay
contacto.

\subsection{Dos rieles de alimentación}
La PCB entrega dos salidas de 3.3\,V independientes desde el 12\,V
constante del OBD2:

\begin{itemize}
  \item \textbf{Riel 1} (siempre encendido): ESP32-S3 + SN65HVD230. Tiene
        que quedar siempre alimentado porque el transceptor es lo único
        capaz de "escuchar" el bus CAN para saber cuándo volvió el
        contacto.
  \item \textbf{Riel 2} (conmutado por software): LEDs + OLED, y por
        separado la microSD. Se corta por completo mientras no hay
        contacto.
\end{itemize}

\subsection{Circuito de \textit{power-gating}: MOSFET y pines}
\label{sec:powergate}
Cada riel 2 se controla con un MOSFET de \textbf{canal P, como interruptor
de alto lado} (corta el 3.3\,V que entra al periférico, no su GND). Se
prefiere sobre un N-MOSFET de bajo lado porque con este último el GND del
periférico puede quedar ligeramente flotante mientras conduce, lo que
genera problemas de integridad de señal en buses referenciados a GND como
I2C (OLED) y SPI (microSD).

\begin{table}[H]
\centering
\caption{Pines de \textit{power-gating}, \texttt{include/config.h}.}
\begin{tabular}{@{}lll@{}}
\toprule
Señal & GPIO ESP32-S3 & Controla \\
\midrule
\texttt{PIN\_PWR\_LED\_OLED} & GPIO6 & LEDs + OLED \\
\texttt{PIN\_PWR\_SD}        & GPIO7 & microSD \\
\bottomrule
\end{tabular}
\end{table}

Ambos GPIO son de dominio RTC en el ESP32-S3 (rango GPIO0-21), requisito
para que \texttt{gpio\_hold\_en()} pueda retener su nivel lógico durante el
\textit{deep sleep} \cite{espressif2024trm}. Topología de cada MOSFET
(misma para los dos rieles):

\begin{figure}[H]
\begin{lstlisting}[style=plain,caption={Interruptor de alto lado con P-MOSFET, un canal por riel.}]
Riel 2 (3.3V) ---+--------------------- Source (P-MOSFET)
                 |
              [R_pullup ~10k]
                 |
GPIO --[R_gate ~1k]------------------- Gate
                                          |
                                        Drain --- carga (OLED+LED, o microSD)
\end{lstlisting}
\end{figure}

El criterio detrás de cada componente: como el GPIO y el riel 2 comparten
el mismo nivel lógico (3.3\,V), no hace falta un transistor adicional para
manejar la compuerta —GPIO en alto da $V_{gs}=0$ (MOSFET apagado); GPIO en
bajo da $V_{gs}\approx-3.3\,\text{V}$, muy por debajo del umbral típico de
un P-MOSFET de nivel lógico ($-1$ a $-2\,\text{V}$ aprox., p.\,ej. un
AO3401A o equivalente)—. La resistencia de \textit{pull-up} en la
compuerta asegura que el MOSFET arranque apagado por defecto durante el
breve instante antes de que el firmware configure el GPIO como salida; la
resistencia en serie hacia el GPIO limita la corriente de carga de la
compuerta.

\subsection{Cadena de regulación}
\label{sec:regchain}
La alimentación completa, del 12\,V del OBD2 a los dos rieles de 3.3\,V,
pasa por tres reguladores en cascada:

\begin{enumerate}
  \item \textbf{TPS54302}: convertidor \textit{buck} (conmutado) que baja
        el 12\,V del pin 16 a 5\,V. Es la única etapa de conversión
        conmutada del sistema; las dos siguientes son lineales.
  \item \textbf{AMS1117} (5\,V $\rightarrow$ 3.3\,V): alimenta el
        \textbf{riel 1} (ESP32-S3 + SN65HVD230). Nunca se apaga.
  \item \textbf{AP2112} (5\,V $\rightarrow$ 3.3\,V): alimenta el
        \textbf{riel 2}; su salida de 3.3\,V entra a los dos MOSFET de la
        Sección~\ref{sec:powergate}, que gatean OLED+LED y microSD por
        separado. El AP2112 en sí mismo nunca se apaga —solo se corta lo
        que hay \textit{después} de los MOSFET—, pero al ser una LDO de
        muy baja corriente de reposo su aporte al consumo en espera es
        pequeño de todos modos (Sección~\ref{sec:consumo}).
\end{enumerate}

Esta cadena (\textit{buck} + 2 LDO en paralelo desde el mismo 5\,V) es una
elección de diseño válida y común, pero tiene una consecuencia directa para
el consumo "sin contacto" que conviene tener presente: a diferencia del
AP2112, el \textbf{AMS1117 es una LDO de propósito general, no de bajo
consumo} —su corriente de reposo (\textit{ground current}) es
comparativamente alta, típicamente 5-8\,mA casi sin importar la carga en
su salida, con variación notable entre fabricantes clon del mismo chip—.
Como el AMS1117 alimenta justo el riel que nunca se apaga, esa corriente de
reposo se sostiene las 24\,horas del día, se detalla en la
Sección~\ref{sec:consumo}.

\subsection{Implementación en firmware}
\texttt{canObd2CheckContact()} manda un único \textit{ping} (PID
\texttt{0x0C}, RPM, obligatorio en todo ECU OBD2) y reporta si algo
respondió, sin importar el valor. Se usa en dos puntos:

\begin{itemize}
  \item Al arrancar (\texttt{setup()}, tanto en un encendido normal como al
        despertar del \textit{deep sleep}): si no hay contacto, el chip
        vuelve a dormir antes de tocar el riel 2 o el WiFi.
  \item Ya despierto, dentro de \texttt{canObd2Task}: si el bus deja de
        responder durante \texttt{CONTACT\_LOST\_DEBOUNCE\_CYCLES} ciclos
        seguidos (3, para no dormirse por un solo \textit{frame} perdido en
        un bus ocupado), se considera que se fue el contacto.
\end{itemize}

\begin{figure}[H]
\begin{lstlisting}[caption={Entrada a deep sleep, \texttt{src/power\_sleep.cpp}.}, label={fig:deepsleep}]
void enterDeepSleepUntilContact() {
  sdLoggerShutdown();  // cierra el archivo antes de cortarle la energia
  WiFi.mode(WIFI_OFF);

  powerGateSet(false);
  gpio_hold_en((gpio_num_t)PIN_PWR_LED_OLED);
  gpio_hold_en((gpio_num_t)PIN_PWR_SD);
  gpio_deep_sleep_hold_en(); // sin esto el hold no sobrevive al deep sleep

  esp_sleep_enable_timer_wakeup(DEEP_SLEEP_WAKE_INTERVAL_US);
  esp_deep_sleep_start();  // no retorna
}
\end{lstlisting}
\end{figure}

Cerrar el archivo de la SD \textbf{antes} de cortar su MOSFET es
obligatorio, no una formalidad: cortar la alimentación con un archivo
abierto arriesga corromper el sistema de archivos completo, no solo perder
la última fila. El \textit{deep sleep} borra toda la RAM salvo el dominio
RTC, así que al despertar \texttt{setup()} vuelve a correr desde el
principio, exactamente como un arranque en frío.

\subsection{Consumo estimado y autonomía de la batería}
\label{sec:consumo}
Los siguientes números son estimaciones a partir de valores típicos de
hoja de datos, no mediciones sobre la placa real —sirven para tener una
idea de orden de magnitud, no un valor exacto—.

\begin{table}[H]
\centering
\small
\caption{Contribuyentes al consumo promedio "sin contacto", cadena real (TPS54302 + AMS1117 + AP2112).}
\begin{tabular}{@{}p{1.7cm}p{3.6cm}p{1.7cm}p{4.8cm}@{}}
\toprule
Nodo & Fuente & Aporte típico & Nota \\
\midrule
3.3\,V (riel 1) & ESP32-S3 en deep sleep & $\sim$0.02\,mA & 10-25\,$\mu$A \cite{espressif2024datasheet} \\
3.3\,V (riel 1) & SN65HVD230, siempre activo & $\sim$5-10\,mA & Modo normal, no standby \cite{sn65hvd230_datasheet} \\
3.3\,V (riel 1) & Pulsos de despertar (cada 20\,s) & $\sim$1.5-2\,mA & Promediado: $\sim$150\,mA por $\sim$0.2-0.3\,s \\
5\,V & \textbf{AMS1117}, corriente de reposo propia & \textbf{$\sim$5-8\,mA} & Casi independiente de la carga \\
5\,V & AP2112, corriente de reposo propia & $\sim$0.06\,mA & Sin carga (MOSFET de salida abiertos) \\
12\,V & TPS54302, corriente de reposo propia & $\sim$0.1-0.3\,mA & \\
\midrule
\multicolumn{2}{@{}l}{\textbf{Total aproximado, lado 12\,V}} & \textbf{$\sim$7-9\,mA} & Ver conversión abajo \\
\bottomrule
\end{tabular}
\end{table}

El resultado es contraintuitivo pero importante: \textbf{el transceptor CAN
y la LDO del riel 1, no el ESP32-S3, dominan el consumo mientras el
dispositivo "duerme"} —porque ese riel no se puede cortar sin perder la
capacidad de detectar cuándo vuelve el contacto (Sección~\ref{sec:powergate}).
El \textit{deep sleep} resuelve la parte del ESP32-S3 casi por completo,
pero no puede resolver la del transceptor ni la de su LDO con esta
arquitectura de dos rieles fija.

Convertir la corriente del lado 3.3\,V/5\,V al lado 12\,V: el AMS1117 y el
AP2112 son LDO (no reducen corriente, solo voltaje, disipando la diferencia
como calor), así que la corriente que el TPS54302 debe entregar en su
salida de 5\,V es aproximadamente la suma de lo que consumen el riel 1
completo ($\sim$8.5\,mA de ESP32+SN65+pulsos) más la corriente de reposo
propia del AMS1117 ($\sim$6.5\,mA punto medio) y del AP2112
($\sim$0.06\,mA): $\approx$15\,mA a 5\,V. El TPS54302, al ser un
\textit{buck} conmutado, sí convierte potencia de forma más eficiente
($\sim$80\,\% en carga liviana, estimado):

\begin{equation}
I_{12V} \approx \frac{5\,\text{V} \times 15\,\text{mA}}{12\,\text{V} \times 0.80} + 0.2\,\text{mA (I}_q\text{ TPS54302)} \approx 8\,\text{mA}
\end{equation}

Tomando esos $\sim$8\,mA y una batería típica de auto de 50\,Ah (de la cual,
por buena práctica, no conviene usar más del 50\,\% para no acortar su vida
útil, es decir $\sim$25\,000\,mAh utilizables):

\begin{equation}
\frac{25\,000\,\text{mAh}}{8\,\text{mA}} \approx 3\,125\,\text{horas} \approx 130\,\text{días}
\end{equation}

Es decir, \textbf{el dispositivo por sí solo} tardaría del orden de 4-5
meses en consumir la mitad de una batería sana —menos que si el riel 1
usara una LDO de baja corriente de reposo en vez del AMS1117, pero todavía
un margen amplio frente al uso normal del vehículo—. En la práctica esto
importa menos de lo que parece, porque el propio vehículo ya tiene su
propio consumo parásito de fábrica (alarma, memorias de radio/ECU, reloj,
etc.), típicamente $\sim$20-50\,mA con todo "apagado" —del mismo orden o
mayor que el aporte de este dispositivo—, y ese consumo de base es el que
realmente determina cuánto aguanta la batería estacionada, con o sin el
ESP32-S3 conectado.

\textbf{Decisión de diseño:} se evaluó reemplazar el AMS1117 del riel 1 por
una LDO de baja corriente de reposo (el propio AP2112, un MCP1700, o una
TPS7A02 con $I_q$ en el orden de nanoamperios bajarían el consumo "sin
contacto" varios mA más), pero se optó por mantener el AMS1117 y el AP2112
tal como estaban ya elegidos en la placa —es el contribuyente individual
más grande después del SN65HVD230, y queda documentado aquí como una
compensación consciente entre simplicidad/costo de la lista de materiales y
la autonomía teórica de la batería, no un descuido.

\subsection{Recomendación de mantenimiento}
Con esos números, la recomendación práctica no cambia mucho respecto al
cuidado normal de cualquier batería de auto: \textbf{arrancar/circular el
vehículo al menos cada 2-3 semanas}, o usar un cargador de mantenimiento
(\textit{battery tender}) si va a quedar estacionado más tiempo que eso.
Con la cadena de reguladores actual (AMS1117 incluido) el dispositivo
agrega una fracción del consumo parásito que el auto ya tiene de fábrica
—notable, pero no dominante—, así que este intervalo sigue siendo una
recomendación razonable sin necesitar ser más estricta.

\subsection{¿Afecta consultar tan seguido?}
Con la arquitectura actual hay dos frecuencias distintas que vale la pena
separar:

\begin{itemize}
  \item \textbf{Sondeo CAN mientras hay contacto} (un ciclo completo de 8
        PIDs cada $\sim$0.2-1.4\,s): el bus CAN a 500\,kbit/s tiene
        capacidad de sobra para esto —cada trama de petición/respuesta son
        unos 100 bits, así que incluso varias por segundo representan bien
        menos del 1\,\% del ancho de banda del bus—. Una ECU real está
        diseñada para atender consultas de diagnóstico bastante más
        frecuentes que esto sin ningún problema; no hay riesgo de saturar
        el bus ni de sobrecargar la ECU.
  \item \textbf{Escritura a la microSD, una fila por segundo}
        (\texttt{SD\_LOG\_INTERVAL\_MS}): las escrituras son secuenciales
        —se van agregando al final del archivo, no se reescribe una y otra
        vez el mismo bloque—, que es el mejor caso posible para el
        desgaste de la memoria flash. Una fila ronda los 100 bytes; a una
        por segundo son unos 360\,KB por hora de viaje, insignificante
        frente a la resistencia típica de una microSD moderna (miles de
        ciclos de escritura por bloque, con \textit{wear leveling}
        integrado). No es un factor de desgaste relevante en la práctica.
\end{itemize}

% ============================================================
\section{Panel web}
\label{sec:panelweb}
% ============================================================

\subsection{Punto de acceso WiFi: por qué SoftAP en vez de solo STA}
El dispositivo siempre levanta su propio punto de acceso (SSID
\texttt{ESP32\_OBD2}, IP típica \texttt{192.168.4.1}) en vez de depender
únicamente de unirse a una red WiFi conocida. El criterio es el contexto de
uso real: dentro de un vehículo casi nunca hay una red WiFi conocida en
rango, así que un diseño solo-STA dejaría el panel inaccesible la mayor
parte del tiempo. El SoftAP garantiza acceso local sin esa dependencia;
unirse además como cliente a una red conocida queda como algo opcional,
únicamente para sincronizar hora por NTP cuando hay una disponible.

\subsection{Por qué ESPAsyncWebServer en vez del servidor síncrono de Arduino}
Se eligió ESPAsyncWebServer \cite{espasyncwebserver2024} sobre el
\texttt{WebServer} síncrono incluido en el core de Arduino-ESP32 porque
atiende las conexiones de forma asíncrona, basada en eventos: una petición
lenta (por ejemplo, una descarga de CSV grande) no bloquea el hilo mientras
espera que el cliente reciba los datos. Dado que el servidor corre en el
mismo núcleo que la pila WiFi (Sección~\ref{sec:herramientas}), un modelo
síncrono habría significado que una descarga lenta retrasara también el
resto del tráfico de red, incluyendo el WebSocket que empuja los datos en
vivo.

\subsection{Empuje en vivo: WebSocket con respaldo por \textit{polling}}
El estado se empuja a los clientes por WebSocket cada 500\,ms
—Figura~\ref{fig:merge} ya mostró de dónde sale ese dato—, con un
mecanismo de respaldo: si el navegador detecta que el WebSocket no está
abierto, recurre a pedir \texttt{/api/live} por HTTP cada 1.5\,s. El
criterio es la robustez: un WebSocket puede caerse por muchas razones ajenas
al firmware (el navegador durmiendo la pestaña, un cambio de red en el
celular), y el panel no debería quedar congelado solo porque ese canal se
cortó.

\subsection{JSON con ArduinoJson}
Las respuestas HTTP y los mensajes WebSocket se serializan con la librería
ArduinoJson \cite{arduinojson2024} en vez de concatenar \textit{strings} a
mano. El criterio: construir JSON por concatenación de texto es una fuente
común de errores de escape/formato difíciles de depurar en un
microcontrolador, mientras que ArduinoJson gestiona el documento en un
búfer de tamaño acotado y conocido, algo relevante en un sistema con 320\,KB
de RAM total repartidos entre cinco tareas.

\begin{table}[H]
\centering
\caption{Endpoints servidos por ESPAsyncWebServer (\texttt{src/web\_server.cpp}).}
\begin{tabular}{@{}lll@{}}
\toprule
Método & Ruta & Descripción \\
\midrule
GET  & \texttt{/}              & Panel del dashboard (HTML embebido en el firmware) \\
WS   & \texttt{/ws}             & WebSocket, empuja el estado en vivo cada 500\,ms \\
GET  & \texttt{/api/live}       & Snapshot JSON del estado actual \\
GET  & \texttt{/api/files}      & Lista de archivos en \texttt{/logs} \\
GET  & \texttt{/api/download}   & Descarga un CSV completo \\
GET  & \texttt{/api/chartdata}  & Puntos decimados para graficar \\
POST & \texttt{/api/reset}      & Reinicia el contador de viaje \\
\bottomrule
\end{tabular}
\end{table}

\subsection{Interfaz}
El HTML/CSS/JS está embebido por completo en el firmware, sin dependencias
de CDN externo —criterio directo del contexto de uso: el SoftAP no tiene
salida a internet, así que cualquier librería externa de gráficas
simplemente no cargaría—. El panel tiene dos pestañas (\textit{En vivo} e
\textit{Historial}), tarjetas destacadas con el consumo instantáneo,
litros y costo del viaje, casillas para elegir qué variables graficar, y un
\textit{tooltip} interactivo con línea guía. Cuando hay 2 o más variables
seleccionadas a la vez, cada una se normaliza a su propio rango (0-100\,\%)
para poder comparar formas de curva con unidades muy distintas (RPM 0-7000
frente a L/100km 0-20) sin que una aplaste visualmente a la otra.

% ============================================================
\section{Compilación y programación}
\label{sec:build}
% ============================================================

\subsection{PlatformIO}

\begin{figure}[H]
\begin{lstlisting}[style=plain,caption={Configuración relevante de \texttt{platformio.ini}.}]
board = esp32-s3-devkitc-1
board_build.flash_mode = qio
board_build.arduino.memory_type = qio_opi
board_build.psram_type = opi
board_upload.flash_size = 16MB
\end{lstlisting}
\end{figure}

Para compilar y subir: \texttt{pio run -t upload} y luego
\texttt{pio device monitor -b 115200}.

\subsection{USB nativo frente a UART externo}
El \textit{devkit} de desarrollo solo expone el \textbf{USB nativo} del
ESP32-S3 (el periférico USB-Serial-JTAG integrado, sin un chip puente
CP2102/CH340 aparte). Esto tiene dos implicancias para una PCB propia:

\begin{enumerate}
  \item Programar por UART \textbf{siempre funciona}, sin importar la
        configuración del firmware: \texttt{esptool} habla directo con el
        \textit{bootloader} de ROM del chip, independiente de cualquier
        opción de compilación de la aplicación.
  \item El firmware define \texttt{-DARDUINO\_USB\_CDC\_ON\_BOOT=1}. Con ese
        \textit{flag}, el objeto \texttt{Serial} de Arduino se enruta al
        USB nativo en vez del UART0 físico (GPIO43/44). Si la PCB propia usa
        un adaptador UART externo y no tiene el USB nativo conectado a
        nada, ese \textit{flag} debe quitarse —de lo contrario el monitor
        serie del adaptador UART no mostrará nada, aunque el firmware esté
        corriendo correctamente.
\end{enumerate}

\subsection{Circuito mínimo para programar/arrancar}
Cualquier placa propia con el ESP32-S3 necesita, como mínimo: \texttt{EN}
(reset) con \textit{pull-up} a 3.3\,V y un capacitor de $\sim$1\,$\mu$F a
GND; \texttt{GPIO0} (BOOT) con \textit{pull-up} a 3.3\,V y un botón a GND
para forzar el modo \textit{bootloader} manualmente; y, opcionalmente, un
circuito de auto-reset (2 transistores desde DTR/RTS del adaptador UART
hacia EN y GPIO0) para que PlatformIO entre solo al modo de programación sin
botones.

% ============================================================
\section{Solución de problemas comunes}
% ============================================================

\begin{itemize}
  \item \textbf{La subida se queda en "Connecting..." indefinidamente.}
        Entrar manualmente en modo \textit{bootloader} (mantener BOOT,
        pulsar y soltar RESET/EN, soltar BOOT, y recién ahí lanzar la
        subida). Verificar también que ningún otro programa tenga el
        puerto COM abierto.
  \item \textbf{El monitor serie no muestra nada tras programar.} Revisar
        que \texttt{ARDUINO\_USB\_CDC\_ON\_BOOT} coincida con el hardware
        real (Sección~\ref{sec:build}).
  \item \textbf{\texttt{sdSelectCard(): Select Failed} / \texttt{SD=0} en el
        arranque.} Casi siempre es cableado/alimentación: revisar tarjeta
        insertada, tensión correcta del módulo lector, continuidad de los 4
        pines SPI + GND, formato FAT32, y el \textit{pull-up} de
        10\,k$\Omega$ en CS si se soldó directo sin módulo.
  \item \textbf{El panel web no carga en el celular.} Escribir
        \texttt{http://192.168.4.1} con el esquema explícito —varios
        navegadores intentan \texttt{https://} por defecto y el dispositivo
        no tiene TLS—. Algunos celulares además desconectan una red WiFi
        sin salida a internet a menos que se les indique quedarse
        conectados igual.
\end{itemize}
 (o copia/pega el
% contenido) dentro de tu documento original, en el punto donde va este
% capítulo. Ajusta el nivel de \section por \subsection si en tu
% documento esto ya cuelga de un \chapter{Diseño del software}.
%
% Requisitos que van en el PREAMBULO de tu documento maestro (no aquí):
%
%   \usepackage{amsmath,amssymb}
%   \usepackage{booktabs}
%   \usepackage{array}
%   \usepackage{xcolor}
%   \usepackage{listings}
%   \usepackage{float}
%   \usepackage{lmodern}    % fuentes escalables -- requisito de microtype
%                           % si tu documento usa babel en espanol (que
%                           % carga fuentes EC de mapa de bits sin esto)
%   \usepackage{microtype}  % reduce cajas "overfull" en parrafos con
%                           % identificadores \texttt largos (p.ej.
%                           % FUEL_CALIBRATION_FACTOR); combinar con:
%   \tolerance=1500
%   \emergencystretch=3em
%
%   \definecolor{codebg}{RGB}{245,245,245}
%   \definecolor{codekw}{RGB}{0,90,170}
%   \definecolor{codecom}{RGB}{110,110,110}
%   \definecolor{codestr}{RGB}{160,40,40}
%   \lstdefinestyle{cpp}{
%     language=C++, backgroundcolor=\color{codebg}, basicstyle=\ttfamily\small,
%     keywordstyle=\color{codekw}\bfseries, commentstyle=\color{codecom}\itshape,
%     stringstyle=\color{codestr}, numbers=left, numberstyle=\tiny\color{gray},
%     numbersep=8pt, breaklines=true, showstringspaces=false,
%     frame=single, rulecolor=\color{gray!40}, tabsize=2,
%   }
%   \lstdefinestyle{plain}{
%     backgroundcolor=\color{codebg}, basicstyle=\ttfamily\small,
%     breaklines=true, showstringspaces=false,
%     frame=single, rulecolor=\color{gray!40},
%   }
%   \lstset{style=cpp}
%
% Bibliografía: las fuentes citadas aquí (\cite{...}) están en
% referencias.bib, para fusionar con la bibliografía única de tu
% documento (donde ya tengas \bibliography{...} o \addbibresource{...}).
% \cite{} es el comando base de LaTeX, compatible con bibliografía clásica
% (\bibliographystyle{apalike} da formato "(Autor, Año)" sin paquetes
% extra), con natbib, y con biblatex.
% ============================================================

Esta sección retoma el capítulo de diseño de software y describe, parte por
parte, el firmware que corre en el ESP32-S3: las herramientas con las que se
construyó, cómo reparte el trabajo entre tareas, el protocolo con el que
habla con el vehículo, el método de cálculo del consumo, cómo persiste los
datos, cómo expone un panel de control, y el criterio detrás de cada
elección de componente y de librería.

% ============================================================
\section{Herramientas y entorno de desarrollo}
\label{sec:herramientas}
% ============================================================

\subsection{Visual Studio Code + PlatformIO}
El firmware se desarrolla en Visual Studio Code con la extensión PlatformIO
IDE, en vez del IDE oficial de Arduino. La diferencia no es cosmética:

\begin{itemize}
  \item \textbf{Toolchain embebido integrado.} PlatformIO gestiona la
        descarga y versión exacta del framework
        (\texttt{framework-arduinoespressif32}) \cite{arduinoesp32core2024},
        el compilador cruzado (\texttt{toolchain-xtensa-esp32s3}) y
        \texttt{esptool} para programar, todo declarado y fijado en
        \texttt{platformio.ini} \cite{platformio2024}. Esto hace que el
        \textit{build} sea reproducible entre máquinas —un requisito real
        una vez que el proyecto tiene múltiples archivos fuente y
        dependencias externas—, algo que el modelo de "sketch único" del
        IDE de Arduino maneja peor.
  \item \textbf{Autocompletado real de C++.} Con un servidor de lenguaje
        indexando todo el árbol de \texttt{src/}/\texttt{include/}, moverse
        entre los siete archivos fuente del proyecto y sus cabeceras
        compartidas (\texttt{config.h}, \texttt{types.h}) es inmediato.
  \item \textbf{Terminal integrado.} Compilar, subir y abrir el monitor
        serie desde el mismo terminal donde se edita el código evita el
        cambio de contexto entre editor y herramienta externa, y deja el
        registro de errores de compilación con enlaces directos a archivo y
        línea.
\end{itemize}

\subsection{Por qué FreeRTOS}
Un sistema que a la vez sondea un bus CAN sensible a tiempos, refresca una
pantalla, escribe en una tarjeta SD y sirve un panel web por WiFi es, por
naturaleza, un problema \textbf{concurrente}. Hay dos formas razonables de
resolverlo: un único \texttt{loop()} con máquinas de estado manuales y
sondeo no bloqueante en cada punto donde se podría esperar —viable en
proyectos pequeños, pero frágil a medida que crecen—, o un sistema operativo
de tiempo real con prioridades y primitivas de sincronización reales. Este
proyecto usa la segunda opción.

FreeRTOS no es, en rigor, una dependencia externa que se "agregó" —es el
planificador que el propio núcleo Arduino-ESP32 ya usa por debajo
\cite{freertos2024}; usarlo explícitamente es aprovechar la plataforma en
vez de pelear contra ella con un \texttt{loop()} monolítico encima.
Concretamente, esto permite:

\begin{itemize}
  \item Darle al sondeo CAN la \textbf{prioridad más alta} de todas las
        tareas, garantizando que una transacción I2C lenta hacia el OLED o
        una descarga de CSV por WiFi nunca retrase la lectura del vehículo.
  \item Proteger el estado compartido con mutex reales en vez de banderas
        improvisadas.
  \item Que agregar una funcionalidad nueva sea, casi siempre, agregar una
        tarea nueva con una prioridad clara, en vez de insertar más lógica
        condicional dentro de un bucle ya complejo —el modo de bajo consumo
        de la Sección~\ref{sec:idle} solo requirió tocar el bucle de la
        tarea CAN y agregar una bandera de estado que las demás tareas ya
        sabían consultar.
\end{itemize}

\begin{figure}[H]
\begin{lstlisting}[caption={Creación de tareas FreeRTOS en \texttt{src/main.cpp}.}, label={fig:tasks}]
statusLedsInit();
xTaskCreatePinnedToCore(statusLedsTask, "leds", STACK_LED_TASK, nullptr,
                         PRIO_LED_TASK, nullptr, CORE_LED_TASK);

bool oledOk = displayInit();
bool sdOk   = sdLoggerInit();
bool canOk  = canObd2Init();

webServerInit(); // reporta ErrorCode::WIFI_FAIL si el softAP falla

if (canOk) {
  xTaskCreatePinnedToCore(canObd2Task, "canObd2", STACK_CAN_TASK, nullptr,
                           PRIO_CAN_TASK, nullptr, CORE_CAN_TASK);
}
if (oledOk) {
  xTaskCreatePinnedToCore(displayTask, "display", STACK_DISPLAY_TASK, nullptr,
                           PRIO_DISPLAY_TASK, nullptr, CORE_DISPLAY_TASK);
}
if (sdOk) {
  xTaskCreatePinnedToCore(sdLoggerTask, "sdLogger", STACK_SD_TASK, nullptr,
                           PRIO_SD_TASK, nullptr, CORE_SD_TASK);
}
\end{lstlisting}
\end{figure}

% ============================================================
\section{Interfaz con el hardware}
% ============================================================

El software está organizado alrededor de las interfaces físicas del
ESP32-S3 \cite{espressif2024datasheet,espressif2024trm}. Cada periférico se
eligió por una razón concreta, no por ser "lo que venía en el kit":

\begin{table}[H]
\centering
\caption{Mapa de pines definido en \texttt{include/config.h}.}
\begin{tabular}{@{}lll@{}}
\toprule
Bus & Señal & GPIO ESP32-S3 \\
\midrule
TWAI (CAN)   & TX $\rightarrow$ SN65HVD230 CTX & GPIO4 \\
TWAI (CAN)   & RX $\leftarrow$ SN65HVD230 CRX  & GPIO5 \\
I2C          & SDA $\rightarrow$ OLED          & GPIO8 \\
I2C          & SCL $\rightarrow$ OLED          & GPIO9 \\
SPI (SD)     & SCK                              & GPIO12 \\
SPI (SD)     & MISO                             & GPIO13 \\
SPI (SD)     & MOSI                             & GPIO11 \\
SPI (SD)     & CS                               & GPIO10 \\
LED          & CAN activo                       & GPIO15 \\
LED          & Actividad SD                     & GPIO16 \\
LED          & Servidor activo                  & GPIO17 \\
LED          & Error                            & GPIO18 \\
\bottomrule
\end{tabular}
\end{table}

\textit{Nota.} Pines reservados —no usar—: GPIO 26-37 (bus interno hacia el
flash y la PSRAM octal del módulo N16R8) y GPIO 0/3/19/20/43-46 (pines de
\textit{strapping}, USB nativo y consola UART0 del chip).

\subsection{Controlador CAN: TWAI integrado, no un chip externo}
El ESP32-S3 trae un controlador CAN (TWAI) integrado en el propio chip
\cite{espressif2024trm}. Esa es la razón por la que el diseño no usa un
controlador CAN externo por SPI como el MCP2515, común en proyectos sobre
Arduino Uno/Nano que no tienen esta capacidad de fábrica: usar el
controlador nativo ahorra un dispositivo, una interrupción y un bus SPI que
ya está ocupado por la microSD, y evita la latencia adicional de tener que
serializar cada trama CAN a través de otro periférico intermediario.

\subsection{Transceptor CAN: SN65HVD230}
El TWAI del ESP32-S3 solo genera niveles lógicos TTL; convertirlos a las
señales diferenciales CAN\_H/CAN\_L del bus físico requiere un transceptor.
Se eligió el SN65HVD230 \cite{sn65hvd230_datasheet} en concreto porque opera
nativamente a 3.3\,V —el mismo nivel lógico del ESP32-S3—, evitando un
\textit{level-shifter} que sí haría falta con transceptores clásicos de 5\,V
como el MCP2551. CAN\_H y CAN\_L se conectan a los pines 6 y 14 del conector
OBD2 (J1962).

\subsection{Pantalla: SSD1306 OLED por I2C}
Se eligió un OLED monocromático de $128\times64$ px con controlador SSD1306
\cite{solomon2008ssd1306} sobre I2C en vez de, por ejemplo, una pantalla TFT
a color por SPI, por tres razones: consume solo 2 pines (I2C) dejando el SPI
libre para la SD; su consumo (10-20\,mA encendida, y se puede apagar por
software, Sección~\ref{sec:idle}) es bajo comparado con un TFT retroiluminado;
y la resolución alcanza de sobra para los pocos números grandes que hace
falta mostrar al conductor —no se necesita color ni gráficos complejos en el
tablero físico, eso ya lo cubre el panel web.

\subsection{Almacenamiento: microSD por SPI dedicado}
La tarjeta microSD usa su propio bus SPI, sin compartirlo con ningún otro
periférico. La alternativa —reutilizar un bus SPI ya en uso— se descartó
porque el registro en la SD ocurre una vez por segundo de forma predecible
mientras que una eventual escritura de mayor tamaño (por ejemplo, servir una
descarga de CSV completa) no debe competir por el bus con otra
responsabilidad no relacionada; un bus dedicado simplifica el análisis de
tiempos y evita tener que coordinar dos periféricos distintos sobre el mismo
SPI.

% ============================================================
\section{Arquitectura de tareas}
% ============================================================

\begin{table}[H]
\centering
\caption{Tareas FreeRTOS, definidas en \texttt{include/config.h} y creadas en \texttt{src/main.cpp}.}
\begin{tabular}{@{}llccc@{}}
\toprule
Tarea & Función & Núcleo & Prioridad & Pila \\
\midrule
CAN/OBD2          & \texttt{canObd2Task}       & 0 & 5 & 4096 \\
SD logger         & \texttt{sdLoggerTask}      & 1 & 4 & 6144 \\
WebSocket push     & \texttt{webSocketPushTask} & 0 & 3 & 4096 \\
OLED display      & \texttt{displayTask}       & 1 & 2 & 4096 \\
LEDs de estado    & \texttt{statusLedsTask}    & 1 & 1 & 2048 \\
\bottomrule
\end{tabular}
\end{table}

El núcleo 0 concentra lo sensible a tiempos (sondeo CAN de alta prioridad +
la pila WiFi/AsyncTCP, que en Arduino-ESP32 corre en ese mismo núcleo). El
núcleo 1 queda con las tareas de interfaz/almacenamiento, que toleran más
\textit{jitter} sin afectar la lectura del vehículo. Las prioridades no son
arbitrarias: reflejan directamente qué tan caro es que cada tarea llegue
tarde (perder un ciclo de sondeo CAN corrompe el cálculo de consumo; llegar
tarde a refrescar un LED es imperceptible).

\subsection{Sincronización entre tareas}
Dos mutex protegen el estado compartido: \texttt{g\_dataMutex} protege
\texttt{g\_vehicleData} (la tarea CAN arma cada ciclo en una copia local y
solo toca la estructura global, bajo mutex, en el paso final de fusión —
Figura~\ref{fig:merge}); \texttt{g\_sdMutex} protege el bus SPI de la
microSD, compartido entre el \textit{logger} periódico y los endpoints de
descarga/gráfica del servidor web. El criterio para usar mutex en vez de,
por ejemplo, deshabilitar interrupciones o copiar datos sin protección, es
que ambos recursos (la estructura de datos y el bus SPI) tienen escritores y
lectores en tareas de distinta prioridad corriendo en núcleos distintos —un
mutex es la primitiva correcta quando el acceso puede, en principio,
bloquear brevemente sin romper los tiempos críticos del sistema.

\begin{figure}[H]
\begin{lstlisting}[caption={Fusión bajo mutex al final de cada ciclo, \texttt{src/can\_obd2.cpp}.}, label={fig:merge}]
if (xSemaphoreTake(g_dataMutex, pdMS_TO_TICKS(100)) == pdTRUE) {
  VehicleData &g = g_vehicleData;
  g.rpm = scratch.rpm;
  g.speed_kmh = scratch.speed_kmh;
  // ... resto de los campos ...
  g.can_status = status;
  g.data_valid = criticalOk;
  g.last_update_ms = now;

  if (criticalOk && dt > 0.0f && dt < 5.0f) {
    fuelCalcUpdate(g, dt);
  }
  xSemaphoreGive(g_dataMutex);
}
\end{lstlisting}
\end{figure}

\subsection{Manejo de errores}
\texttt{systemReportError(ErrorCode)} centraliza el reporte de fallos: guarda
el último código en \texttt{g\_systemStatus.error\_code} (dispara el parpadeo
del LED de error) y lo imprime por \texttt{Serial} con el detalle completo.
El criterio es separar dos audiencias distintas: el LED es para el
conductor, que solo necesita saber \textit{que algo falló}; el log serie es
para depuración en banco, donde sí importa el detalle completo y el
historial.

\begin{table}[H]
\centering
\caption{Códigos de \texttt{ErrorCode} (\texttt{include/types.h}).}
\begin{tabular}{@{}cl@{}}
\toprule
Código & Significado \\
\midrule
1 & Fallo al iniciar el driver TWAI/CAN \\
2 & Fallo al iniciar la microSD \\
3 & Fallo al iniciar el OLED \\
4 & Fallo al iniciar el punto de acceso WiFi \\
\bottomrule
\end{tabular}
\end{table}

% ============================================================
\section{Protocolo CAN / OBD2}
% ============================================================

El controlador TWAI del ESP32-S3 se configura a \textbf{500\,kbit/s}, el
estándar de la mayoría de vehículos livianos con OBD2 sobre CAN
(ISO 15765-4, \cite{iso157654}). Todas las
peticiones usadas caben en un único \textit{frame} CAN de 8 bytes
(\textit{single frame}), así que no hace falta implementar el mecanismo
multi-\textit{frame} de control de flujo ISO-TP —una simplificación
deliberada: ninguno de los PIDs que este proyecto necesita supera ese
tamaño, y soportar multi-\textit{frame} sin necesitarlo solo agrega
superficie de error.

\subsection{PIDs consultados (Modo 01)}
El conjunto de PIDs se eligió por ser el mínimo necesario para el cálculo
speed-density (Sección~\ref{sec:speeddensity}) más los que enriquecen el
registro sin costo relevante de tiempo de bus, siguiendo el estándar SAE
J1979 \cite{sae_j1979_2014}:

\begin{table}[H]
\centering
\caption{PIDs Mode 01 consultados por \texttt{canObd2Task}.}
\begin{tabular}{@{}cll@{}}
\toprule
PID & Dato & Uso \\
\midrule
\texttt{0x0C} & RPM                          & speed-density (crítico) \\
\texttt{0x0B} & MAP (kPa)                    & speed-density (crítico) \\
\texttt{0x0F} & IAT (\textdegree C)          & speed-density (crítico) \\
\texttt{0x0D} & Velocidad (km/h)             & L/100km, distancia (crítico) \\
\texttt{0x05} & Temperatura refrigerante     & mostrado en pantalla/CSV \\
\texttt{0x04} & Carga calculada (\%)         & mostrado en pantalla/CSV \\
\texttt{0x06} & Fuel trim corto plazo (\%)   & corrige el AFR estimado \\
\texttt{0x07} & Fuel trim largo plazo (\%)   & corrige el AFR estimado \\
\bottomrule
\end{tabular}
\end{table}

\textit{Nota.} Los cuatro PIDs marcados "crítico" son indispensables para el
cálculo speed-density; si la ECU no responde alguno, el ciclo completo se
marca inválido y no se integra al viaje.

\subsection{Formato de la petición}

\begin{figure}[H]
\begin{lstlisting}[style=plain,caption={Trama de petición Mode 01, PID genérico.}]
ID = 0x7DF
Data[0] = 0x02   ; SF, 2 bytes de payload (modo + pid)
Data[1] = 0x01   ; Modo 01: datos actuales
Data[2] = <PID>
Data[3..7] = 0x00 (relleno)
\end{lstlisting}
\end{figure}

Las ECUs responden en el rango \texttt{0x7E8}-\texttt{0x7EF}; se acepta la
primera respuesta cuyo \texttt{Data[1] == 0x41} y \texttt{Data[2]} coincida
con el PID pedido.

\subsection{Recuperación de bus-off}
Si el controlador entra en estado \texttt{BUS\_OFF}, la recuperación se
sondea en cada vuelta del bucle sin bloquear con una espera fija —según el
driver TWAI de ESP-IDF \cite{espressif2024trm}, la recuperación solo llega a
\texttt{STOPPED} tras contar 128 ocurrencias de 11 bits recesivos
consecutivos en el bus, un tiempo que depende del tráfico real y no se puede
predecir de antemano; sondear en vez de esperar un tiempo fijo evita tanto
reintentar demasiado pronto (fallar en vano) como esperar de más
innecesariamente.

% ============================================================
\section{Método de cálculo: Speed-Density}
\label{sec:speeddensity}
% ============================================================

El método speed-density estima el flujo másico de aire admitido por el motor
a partir de la ley de gases ideales, sin medirlo con un sensor MAF —la
razón de fondo para elegir este método sobre uno basado en MAF real es que
la mayoría de los vehículos con OBD2 reportan MAP, no flujo de aire directo,
así que speed-density es el único método viable sin instalar un sensor
adicional—:

\begin{equation}
\dot{m}_{aire} \;[\text{g/s}] =
\frac{\dfrac{VE}{100} \cdot RPM \cdot MAP\,[\text{kPa}] \cdot V_d\,[\text{L}] \cdot 28.97}
     {2 \cdot 8.314 \cdot T_{IAT}\,[\text{K}] \cdot 60}
\end{equation}

donde $VE$ es la eficiencia volumétrica (\%) interpolada de una tabla
calibrable (Sección~\ref{sec:ve}); $V_d$ la cilindrada en litros; $28.97$ la
masa molar del aire (g/mol); $8.314$ la constante universal de los gases
(L$\cdot$kPa/(mol$\cdot$K)); el factor $2$ corresponde a que un motor de 4
tiempos completa una admisión cada 2 vueltas de cigüeñal; y $T_{IAT}$ la
temperatura de admisión en Kelvin, con guarda de rango
[$-40$,\,$100$]\,\textdegree C.

\begin{figure}[H]
\begin{lstlisting}[caption={Núcleo del cálculo speed-density, \texttt{src/fuel\_calc.cpp}.}, label={fig:fuelcalc}]
float ve = fuelCalcVeForRpm(vd.rpm);
float iat_K = vd.iat_c + 273.15f;
if (iat_K < 233.15f || iat_K > 373.15f) iat_K = 288.15f;

const float MM_AIR = 28.97f; // g/mol
const float R = 8.314f;      // L*kPa/(mol*K)

float maf = (ve / 100.0f) * (float)vd.rpm * vd.map_kpa *
            ENGINE_DISPLACEMENT_L * MM_AIR /
            (2.0f * R * iat_K * 60.0f);
\end{lstlisting}
\end{figure}

\subsection{De aire a combustible}

\begin{equation}
\dot{m}_{comb}\,[\text{g/s}] = \frac{\dot{m}_{aire}}{AFR} \cdot k_{cal}
\end{equation}

$AFR$ es la relación aire/combustible, por defecto el valor estequiométrico
de la gasolina ($AFR_{stoich} = 14.7$), opcionalmente corregido con los
\textit{fuel trims} que reporta la ECU:

\begin{equation}
AFR_{efectivo} = \frac{AFR_{stoich}}{1 + \frac{STFT\% + LTFT\%}{100}}
\end{equation}

acotado a $[8, 25]$. $k_{cal}$ es \texttt{FUEL\_CALIBRATION\_FACTOR}
(Sección~\ref{sec:calibracion}), $1.0$ por defecto. El criterio para incluir
la corrección por \textit{fuel trim} es que, sin ella, el cálculo asume que
el motor siempre corre exactamente en estequiométrico, cosa que en la
práctica no ocurre; usar los \textit{trims} que la propia ECU ya reporta es
una corrección "gratis" —no requiere sensores adicionales— frente a un error
sistemático conocido.

\subsection{De masa a volumen}

\begin{align}
\dot{V}_{comb}\,[\text{L/h}] &= \frac{\dot{m}_{comb} \cdot 3600}{\rho_{comb}} \\[4pt]
L/100km &=
  \begin{cases}
    \dfrac{\dot{V}_{comb}}{v}\cdot 100 & \text{si } v > 2\,\text{km/h} \\[6pt]
    0 & \text{vehículo detenido}
  \end{cases}
\end{align}

con $\rho_{comb} = 745\,\text{g/L}$ y $v$ la velocidad en km/h.

\subsection{Integración del viaje}

\begin{align}
L_{viaje} &\mathrel{+}= \dot{m}_{comb} \cdot dt \,/\, \rho_{comb} \\
km_{viaje} &\mathrel{+}= v \cdot dt / 3600 \\
Bs_{viaje} &= L_{viaje} \cdot P_{L}
\end{align}

con $dt$ el paso de tiempo real entre ciclos —no un valor fijo, precisamente
para no perder precisión si un ciclo de sondeo tarda más o menos de lo
típico— y $P_L = 6.96$ (\texttt{FUEL\_PRICE\_BS\_PER\_L}) el precio de
referencia del litro.

\subsection{Tabla de eficiencia volumétrica (VE)}
\label{sec:ve}

\begin{table}[H]
\centering
\caption{\texttt{kVeTable}, \texttt{src/fuel\_calc.cpp}. Interpolación lineal entre puntos.}
\begin{tabular}{@{}cc@{}}
\toprule
RPM & VE (\%) \\
\midrule
700  & 60.0 \\
1500 & 74.0 \\
2500 & 85.0 \\
4000 & 88.0 \\
5500 & 83.0 \\
7000 & 74.0 \\
\bottomrule
\end{tabular}
\end{table}

\textit{Nota.} Esta tabla es un punto de partida genérico, no un dato de
fábrica del vehículo. Sin calibrarla, el consumo calculado tendrá un error
sistemático.

\subsection{Calibración contra un tanque real}
\label{sec:calibracion}
No existe forma de que el ESP32 sepa, por sí solo, si el número que calcula
es correcto —no hay con qué comparar dentro del propio microcontrolador—.
El procedimiento recomendado: llenar el tanque y anotar el odómetro, manejar
con normalidad hasta la próxima recarga, comparar los litros reales cargados
contra \texttt{trip\_fuel\_L} acumulado en esa misma distancia, y ajustar

\begin{equation}
\texttt{FUEL\_CALIBRATION\_FACTOR} = \frac{L_{real}}{L_{calculado}}
\end{equation}

en \texttt{config.h}. El criterio de aplicar un único factor multiplicativo
sobre \texttt{fuel\_gs}, en vez de reajustar toda la tabla VE a mano, es que
corrige de forma consistente el instantáneo, el acumulado del viaje y el
costo con un solo número, y es reversible/auditable —a diferencia de
retocar seis puntos de la tabla VE por prueba y error.

% ============================================================
\section{Registro en microSD}
% ============================================================

Cada arranque crea un archivo nuevo \texttt{/logs/trip\_NNN.csv} (índice
autoincremental calculado escaneando los archivos existentes, en vez de un
contador separado que podría desincronizarse tras un corte de energía). Se
escribe una fila por segundo, solo después del primer ciclo CAN válido, y se
pausa por completo durante el modo de bajo consumo (Sección~\ref{sec:idle}).

\subsection{Por qué CSV}
Se eligió CSV en vez de un formato binario compacto o JSON por fila por tres
razones: es directamente legible/abrible en cualquier hoja de cálculo sin
escribir un parser; el tamaño de fila (unos 100 bytes) es intrascendente
frente a la capacidad de una microSD moderna incluso a una fila por segundo
durante horas; y es el formato que el propio panel web ya necesita poder
descargar y volver a interpretar (Sección~\ref{sec:panelweb}), así que un
mismo formato sirve para ambos usos sin conversión intermedia.

\begin{figure}[H]
\begin{lstlisting}[style=plain,caption={Fila de ejemplo del CSV (encabezado + una fila de datos).}]
uptime_ms,rpm,speed_kmh,map_kpa,iat_c,coolant_c,ve_pct,maf_gs,
fuel_gs,instant_Lh,instant_L100km,trip_km,trip_fuel_L,
trip_avg_L100km,can_status,trip_cost_bs

184532,2150,62.0,58.0,29.0,89.0,81.4,7.82,0.532,1.91,3.09,
12.400,0.384,3.10,1,2.67
\end{lstlisting}
\end{figure}

\begin{table}[H]
\centering
\caption{Columnas escritas por \texttt{sdLoggerTask} (\texttt{src/sd\_logger.cpp}).}
\begin{tabular}{@{}ll@{}}
\toprule
Columna & Descripción \\
\midrule
\texttt{uptime\_ms}        & Milisegundos desde el arranque \\
\texttt{rpm}                & Revoluciones por minuto \\
\texttt{speed\_kmh}        & Velocidad, km/h \\
\texttt{map\_kpa}          & Presión absoluta del múltiple, kPa \\
\texttt{iat\_c}             & Temperatura del aire de admisión, \textdegree C \\
\texttt{coolant\_c}        & Temperatura del refrigerante, \textdegree C \\
\texttt{ve\_pct}            & Eficiencia volumétrica usada, \% \\
\texttt{maf\_gs}            & Flujo másico de aire calculado, g/s \\
\texttt{fuel\_gs}           & Flujo másico de combustible calculado, g/s \\
\texttt{instant\_Lh}       & Consumo instantáneo, L/h \\
\texttt{instant\_L100km}  & Consumo instantáneo, L/100km \\
\texttt{trip\_km}          & Distancia acumulada del viaje, km \\
\texttt{trip\_fuel\_L}    & Combustible acumulado del viaje, L \\
\texttt{trip\_avg\_L100km} & Promedio del viaje, L/100km \\
\texttt{can\_status}       & Estado del bus CAN (enum \texttt{CanStatus}) \\
\texttt{trip\_cost\_bs}    & Costo acumulado del viaje, Bs \\
\bottomrule
\end{tabular}
\end{table}

% ============================================================
\section{Modo de bajo consumo ("sin contacto")}
\label{sec:idle}
% ============================================================

\subsection{Motivación}
El dispositivo se alimenta del pin 16 del conector OBD2, que entrega
12\,V \textbf{constantes} sin importar la posición de la llave. Si el
firmware siguiera sondeando el bus a ritmo completo, con el panel web y
todos los periféricos encendidos las 24 horas, se convertiría en una carga
parásita permanente sobre la batería del vehículo. La primera versión de
este modo se basaba en detectar RPM$=0$ y solo bajaba el ritmo de sondeo
—una mejora real pero modesta, porque el ESP32-S3 seguía activo con el
WiFi encendido—. La versión actual va más lejos: separa la alimentación en
dos rieles y usa \textbf{deep sleep} real del ESP32-S3 mientras no hay
contacto.

\subsection{Dos rieles de alimentación}
La PCB entrega dos salidas de 3.3\,V independientes desde el 12\,V
constante del OBD2:

\begin{itemize}
  \item \textbf{Riel 1} (siempre encendido): ESP32-S3 + SN65HVD230. Tiene
        que quedar siempre alimentado porque el transceptor es lo único
        capaz de "escuchar" el bus CAN para saber cuándo volvió el
        contacto.
  \item \textbf{Riel 2} (conmutado por software): LEDs + OLED, y por
        separado la microSD. Se corta por completo mientras no hay
        contacto.
\end{itemize}

\subsection{Circuito de \textit{power-gating}: MOSFET y pines}
\label{sec:powergate}
Cada riel 2 se controla con un MOSFET de \textbf{canal P, como interruptor
de alto lado} (corta el 3.3\,V que entra al periférico, no su GND). Se
prefiere sobre un N-MOSFET de bajo lado porque con este último el GND del
periférico puede quedar ligeramente flotante mientras conduce, lo que
genera problemas de integridad de señal en buses referenciados a GND como
I2C (OLED) y SPI (microSD).

\begin{table}[H]
\centering
\caption{Pines de \textit{power-gating}, \texttt{include/config.h}.}
\begin{tabular}{@{}lll@{}}
\toprule
Señal & GPIO ESP32-S3 & Controla \\
\midrule
\texttt{PIN\_PWR\_LED\_OLED} & GPIO6 & LEDs + OLED \\
\texttt{PIN\_PWR\_SD}        & GPIO7 & microSD \\
\bottomrule
\end{tabular}
\end{table}

Ambos GPIO son de dominio RTC en el ESP32-S3 (rango GPIO0-21), requisito
para que \texttt{gpio\_hold\_en()} pueda retener su nivel lógico durante el
\textit{deep sleep} \cite{espressif2024trm}. Topología de cada MOSFET
(misma para los dos rieles):

\begin{figure}[H]
\begin{lstlisting}[style=plain,caption={Interruptor de alto lado con P-MOSFET, un canal por riel.}]
Riel 2 (3.3V) ---+--------------------- Source (P-MOSFET)
                 |
              [R_pullup ~10k]
                 |
GPIO --[R_gate ~1k]------------------- Gate
                                          |
                                        Drain --- carga (OLED+LED, o microSD)
\end{lstlisting}
\end{figure}

El criterio detrás de cada componente: como el GPIO y el riel 2 comparten
el mismo nivel lógico (3.3\,V), no hace falta un transistor adicional para
manejar la compuerta —GPIO en alto da $V_{gs}=0$ (MOSFET apagado); GPIO en
bajo da $V_{gs}\approx-3.3\,\text{V}$, muy por debajo del umbral típico de
un P-MOSFET de nivel lógico ($-1$ a $-2\,\text{V}$ aprox., p.\,ej. un
AO3401A o equivalente)—. La resistencia de \textit{pull-up} en la
compuerta asegura que el MOSFET arranque apagado por defecto durante el
breve instante antes de que el firmware configure el GPIO como salida; la
resistencia en serie hacia el GPIO limita la corriente de carga de la
compuerta.

\subsection{Cadena de regulación}
\label{sec:regchain}
La alimentación completa, del 12\,V del OBD2 a los dos rieles de 3.3\,V,
pasa por tres reguladores en cascada:

\begin{enumerate}
  \item \textbf{TPS54302}: convertidor \textit{buck} (conmutado) que baja
        el 12\,V del pin 16 a 5\,V. Es la única etapa de conversión
        conmutada del sistema; las dos siguientes son lineales.
  \item \textbf{AMS1117} (5\,V $\rightarrow$ 3.3\,V): alimenta el
        \textbf{riel 1} (ESP32-S3 + SN65HVD230). Nunca se apaga.
  \item \textbf{AP2112} (5\,V $\rightarrow$ 3.3\,V): alimenta el
        \textbf{riel 2}; su salida de 3.3\,V entra a los dos MOSFET de la
        Sección~\ref{sec:powergate}, que gatean OLED+LED y microSD por
        separado. El AP2112 en sí mismo nunca se apaga —solo se corta lo
        que hay \textit{después} de los MOSFET—, pero al ser una LDO de
        muy baja corriente de reposo su aporte al consumo en espera es
        pequeño de todos modos (Sección~\ref{sec:consumo}).
\end{enumerate}

Esta cadena (\textit{buck} + 2 LDO en paralelo desde el mismo 5\,V) es una
elección de diseño válida y común, pero tiene una consecuencia directa para
el consumo "sin contacto" que conviene tener presente: a diferencia del
AP2112, el \textbf{AMS1117 es una LDO de propósito general, no de bajo
consumo} —su corriente de reposo (\textit{ground current}) es
comparativamente alta, típicamente 5-8\,mA casi sin importar la carga en
su salida, con variación notable entre fabricantes clon del mismo chip—.
Como el AMS1117 alimenta justo el riel que nunca se apaga, esa corriente de
reposo se sostiene las 24\,horas del día, se detalla en la
Sección~\ref{sec:consumo}.

\subsection{Implementación en firmware}
\texttt{canObd2CheckContact()} manda un único \textit{ping} (PID
\texttt{0x0C}, RPM, obligatorio en todo ECU OBD2) y reporta si algo
respondió, sin importar el valor. Se usa en dos puntos:

\begin{itemize}
  \item Al arrancar (\texttt{setup()}, tanto en un encendido normal como al
        despertar del \textit{deep sleep}): si no hay contacto, el chip
        vuelve a dormir antes de tocar el riel 2 o el WiFi.
  \item Ya despierto, dentro de \texttt{canObd2Task}: si el bus deja de
        responder durante \texttt{CONTACT\_LOST\_DEBOUNCE\_CYCLES} ciclos
        seguidos (3, para no dormirse por un solo \textit{frame} perdido en
        un bus ocupado), se considera que se fue el contacto.
\end{itemize}

\begin{figure}[H]
\begin{lstlisting}[caption={Entrada a deep sleep, \texttt{src/power\_sleep.cpp}.}, label={fig:deepsleep}]
void enterDeepSleepUntilContact() {
  sdLoggerShutdown();  // cierra el archivo antes de cortarle la energia
  WiFi.mode(WIFI_OFF);

  powerGateSet(false);
  gpio_hold_en((gpio_num_t)PIN_PWR_LED_OLED);
  gpio_hold_en((gpio_num_t)PIN_PWR_SD);
  gpio_deep_sleep_hold_en(); // sin esto el hold no sobrevive al deep sleep

  esp_sleep_enable_timer_wakeup(DEEP_SLEEP_WAKE_INTERVAL_US);
  esp_deep_sleep_start();  // no retorna
}
\end{lstlisting}
\end{figure}

Cerrar el archivo de la SD \textbf{antes} de cortar su MOSFET es
obligatorio, no una formalidad: cortar la alimentación con un archivo
abierto arriesga corromper el sistema de archivos completo, no solo perder
la última fila. El \textit{deep sleep} borra toda la RAM salvo el dominio
RTC, así que al despertar \texttt{setup()} vuelve a correr desde el
principio, exactamente como un arranque en frío.

\subsection{Consumo estimado y autonomía de la batería}
\label{sec:consumo}
Los siguientes números son estimaciones a partir de valores típicos de
hoja de datos, no mediciones sobre la placa real —sirven para tener una
idea de orden de magnitud, no un valor exacto—.

\begin{table}[H]
\centering
\small
\caption{Contribuyentes al consumo promedio "sin contacto", cadena real (TPS54302 + AMS1117 + AP2112).}
\begin{tabular}{@{}p{1.7cm}p{3.6cm}p{1.7cm}p{4.8cm}@{}}
\toprule
Nodo & Fuente & Aporte típico & Nota \\
\midrule
3.3\,V (riel 1) & ESP32-S3 en deep sleep & $\sim$0.02\,mA & 10-25\,$\mu$A \cite{espressif2024datasheet} \\
3.3\,V (riel 1) & SN65HVD230, siempre activo & $\sim$5-10\,mA & Modo normal, no standby \cite{sn65hvd230_datasheet} \\
3.3\,V (riel 1) & Pulsos de despertar (cada 20\,s) & $\sim$1.5-2\,mA & Promediado: $\sim$150\,mA por $\sim$0.2-0.3\,s \\
5\,V & \textbf{AMS1117}, corriente de reposo propia & \textbf{$\sim$5-8\,mA} & Casi independiente de la carga \\
5\,V & AP2112, corriente de reposo propia & $\sim$0.06\,mA & Sin carga (MOSFET de salida abiertos) \\
12\,V & TPS54302, corriente de reposo propia & $\sim$0.1-0.3\,mA & \\
\midrule
\multicolumn{2}{@{}l}{\textbf{Total aproximado, lado 12\,V}} & \textbf{$\sim$7-9\,mA} & Ver conversión abajo \\
\bottomrule
\end{tabular}
\end{table}

El resultado es contraintuitivo pero importante: \textbf{el transceptor CAN
y la LDO del riel 1, no el ESP32-S3, dominan el consumo mientras el
dispositivo "duerme"} —porque ese riel no se puede cortar sin perder la
capacidad de detectar cuándo vuelve el contacto (Sección~\ref{sec:powergate}).
El \textit{deep sleep} resuelve la parte del ESP32-S3 casi por completo,
pero no puede resolver la del transceptor ni la de su LDO con esta
arquitectura de dos rieles fija.

Convertir la corriente del lado 3.3\,V/5\,V al lado 12\,V: el AMS1117 y el
AP2112 son LDO (no reducen corriente, solo voltaje, disipando la diferencia
como calor), así que la corriente que el TPS54302 debe entregar en su
salida de 5\,V es aproximadamente la suma de lo que consumen el riel 1
completo ($\sim$8.5\,mA de ESP32+SN65+pulsos) más la corriente de reposo
propia del AMS1117 ($\sim$6.5\,mA punto medio) y del AP2112
($\sim$0.06\,mA): $\approx$15\,mA a 5\,V. El TPS54302, al ser un
\textit{buck} conmutado, sí convierte potencia de forma más eficiente
($\sim$80\,\% en carga liviana, estimado):

\begin{equation}
I_{12V} \approx \frac{5\,\text{V} \times 15\,\text{mA}}{12\,\text{V} \times 0.80} + 0.2\,\text{mA (I}_q\text{ TPS54302)} \approx 8\,\text{mA}
\end{equation}

Tomando esos $\sim$8\,mA y una batería típica de auto de 50\,Ah (de la cual,
por buena práctica, no conviene usar más del 50\,\% para no acortar su vida
útil, es decir $\sim$25\,000\,mAh utilizables):

\begin{equation}
\frac{25\,000\,\text{mAh}}{8\,\text{mA}} \approx 3\,125\,\text{horas} \approx 130\,\text{días}
\end{equation}

Es decir, \textbf{el dispositivo por sí solo} tardaría del orden de 4-5
meses en consumir la mitad de una batería sana —menos que si el riel 1
usara una LDO de baja corriente de reposo en vez del AMS1117, pero todavía
un margen amplio frente al uso normal del vehículo—. En la práctica esto
importa menos de lo que parece, porque el propio vehículo ya tiene su
propio consumo parásito de fábrica (alarma, memorias de radio/ECU, reloj,
etc.), típicamente $\sim$20-50\,mA con todo "apagado" —del mismo orden o
mayor que el aporte de este dispositivo—, y ese consumo de base es el que
realmente determina cuánto aguanta la batería estacionada, con o sin el
ESP32-S3 conectado.

\textbf{Decisión de diseño:} se evaluó reemplazar el AMS1117 del riel 1 por
una LDO de baja corriente de reposo (el propio AP2112, un MCP1700, o una
TPS7A02 con $I_q$ en el orden de nanoamperios bajarían el consumo "sin
contacto" varios mA más), pero se optó por mantener el AMS1117 y el AP2112
tal como estaban ya elegidos en la placa —es el contribuyente individual
más grande después del SN65HVD230, y queda documentado aquí como una
compensación consciente entre simplicidad/costo de la lista de materiales y
la autonomía teórica de la batería, no un descuido.

\subsection{Recomendación de mantenimiento}
Con esos números, la recomendación práctica no cambia mucho respecto al
cuidado normal de cualquier batería de auto: \textbf{arrancar/circular el
vehículo al menos cada 2-3 semanas}, o usar un cargador de mantenimiento
(\textit{battery tender}) si va a quedar estacionado más tiempo que eso.
Con la cadena de reguladores actual (AMS1117 incluido) el dispositivo
agrega una fracción del consumo parásito que el auto ya tiene de fábrica
—notable, pero no dominante—, así que este intervalo sigue siendo una
recomendación razonable sin necesitar ser más estricta.

\subsection{¿Afecta consultar tan seguido?}
Con la arquitectura actual hay dos frecuencias distintas que vale la pena
separar:

\begin{itemize}
  \item \textbf{Sondeo CAN mientras hay contacto} (un ciclo completo de 8
        PIDs cada $\sim$0.2-1.4\,s): el bus CAN a 500\,kbit/s tiene
        capacidad de sobra para esto —cada trama de petición/respuesta son
        unos 100 bits, así que incluso varias por segundo representan bien
        menos del 1\,\% del ancho de banda del bus—. Una ECU real está
        diseñada para atender consultas de diagnóstico bastante más
        frecuentes que esto sin ningún problema; no hay riesgo de saturar
        el bus ni de sobrecargar la ECU.
  \item \textbf{Escritura a la microSD, una fila por segundo}
        (\texttt{SD\_LOG\_INTERVAL\_MS}): las escrituras son secuenciales
        —se van agregando al final del archivo, no se reescribe una y otra
        vez el mismo bloque—, que es el mejor caso posible para el
        desgaste de la memoria flash. Una fila ronda los 100 bytes; a una
        por segundo son unos 360\,KB por hora de viaje, insignificante
        frente a la resistencia típica de una microSD moderna (miles de
        ciclos de escritura por bloque, con \textit{wear leveling}
        integrado). No es un factor de desgaste relevante en la práctica.
\end{itemize}

% ============================================================
\section{Panel web}
\label{sec:panelweb}
% ============================================================

\subsection{Punto de acceso WiFi: por qué SoftAP en vez de solo STA}
El dispositivo siempre levanta su propio punto de acceso (SSID
\texttt{ESP32\_OBD2}, IP típica \texttt{192.168.4.1}) en vez de depender
únicamente de unirse a una red WiFi conocida. El criterio es el contexto de
uso real: dentro de un vehículo casi nunca hay una red WiFi conocida en
rango, así que un diseño solo-STA dejaría el panel inaccesible la mayor
parte del tiempo. El SoftAP garantiza acceso local sin esa dependencia;
unirse además como cliente a una red conocida queda como algo opcional,
únicamente para sincronizar hora por NTP cuando hay una disponible.

\subsection{Por qué ESPAsyncWebServer en vez del servidor síncrono de Arduino}
Se eligió ESPAsyncWebServer \cite{espasyncwebserver2024} sobre el
\texttt{WebServer} síncrono incluido en el core de Arduino-ESP32 porque
atiende las conexiones de forma asíncrona, basada en eventos: una petición
lenta (por ejemplo, una descarga de CSV grande) no bloquea el hilo mientras
espera que el cliente reciba los datos. Dado que el servidor corre en el
mismo núcleo que la pila WiFi (Sección~\ref{sec:herramientas}), un modelo
síncrono habría significado que una descarga lenta retrasara también el
resto del tráfico de red, incluyendo el WebSocket que empuja los datos en
vivo.

\subsection{Empuje en vivo: WebSocket con respaldo por \textit{polling}}
El estado se empuja a los clientes por WebSocket cada 500\,ms
—Figura~\ref{fig:merge} ya mostró de dónde sale ese dato—, con un
mecanismo de respaldo: si el navegador detecta que el WebSocket no está
abierto, recurre a pedir \texttt{/api/live} por HTTP cada 1.5\,s. El
criterio es la robustez: un WebSocket puede caerse por muchas razones ajenas
al firmware (el navegador durmiendo la pestaña, un cambio de red en el
celular), y el panel no debería quedar congelado solo porque ese canal se
cortó.

\subsection{JSON con ArduinoJson}
Las respuestas HTTP y los mensajes WebSocket se serializan con la librería
ArduinoJson \cite{arduinojson2024} en vez de concatenar \textit{strings} a
mano. El criterio: construir JSON por concatenación de texto es una fuente
común de errores de escape/formato difíciles de depurar en un
microcontrolador, mientras que ArduinoJson gestiona el documento en un
búfer de tamaño acotado y conocido, algo relevante en un sistema con 320\,KB
de RAM total repartidos entre cinco tareas.

\begin{table}[H]
\centering
\caption{Endpoints servidos por ESPAsyncWebServer (\texttt{src/web\_server.cpp}).}
\begin{tabular}{@{}lll@{}}
\toprule
Método & Ruta & Descripción \\
\midrule
GET  & \texttt{/}              & Panel del dashboard (HTML embebido en el firmware) \\
WS   & \texttt{/ws}             & WebSocket, empuja el estado en vivo cada 500\,ms \\
GET  & \texttt{/api/live}       & Snapshot JSON del estado actual \\
GET  & \texttt{/api/files}      & Lista de archivos en \texttt{/logs} \\
GET  & \texttt{/api/download}   & Descarga un CSV completo \\
GET  & \texttt{/api/chartdata}  & Puntos decimados para graficar \\
POST & \texttt{/api/reset}      & Reinicia el contador de viaje \\
\bottomrule
\end{tabular}
\end{table}

\subsection{Interfaz}
El HTML/CSS/JS está embebido por completo en el firmware, sin dependencias
de CDN externo —criterio directo del contexto de uso: el SoftAP no tiene
salida a internet, así que cualquier librería externa de gráficas
simplemente no cargaría—. El panel tiene dos pestañas (\textit{En vivo} e
\textit{Historial}), tarjetas destacadas con el consumo instantáneo,
litros y costo del viaje, casillas para elegir qué variables graficar, y un
\textit{tooltip} interactivo con línea guía. Cuando hay 2 o más variables
seleccionadas a la vez, cada una se normaliza a su propio rango (0-100\,\%)
para poder comparar formas de curva con unidades muy distintas (RPM 0-7000
frente a L/100km 0-20) sin que una aplaste visualmente a la otra.

% ============================================================
\section{Compilación y programación}
\label{sec:build}
% ============================================================

\subsection{PlatformIO}

\begin{figure}[H]
\begin{lstlisting}[style=plain,caption={Configuración relevante de \texttt{platformio.ini}.}]
board = esp32-s3-devkitc-1
board_build.flash_mode = qio
board_build.arduino.memory_type = qio_opi
board_build.psram_type = opi
board_upload.flash_size = 16MB
\end{lstlisting}
\end{figure}

Para compilar y subir: \texttt{pio run -t upload} y luego
\texttt{pio device monitor -b 115200}.

\subsection{USB nativo frente a UART externo}
El \textit{devkit} de desarrollo solo expone el \textbf{USB nativo} del
ESP32-S3 (el periférico USB-Serial-JTAG integrado, sin un chip puente
CP2102/CH340 aparte). Esto tiene dos implicancias para una PCB propia:

\begin{enumerate}
  \item Programar por UART \textbf{siempre funciona}, sin importar la
        configuración del firmware: \texttt{esptool} habla directo con el
        \textit{bootloader} de ROM del chip, independiente de cualquier
        opción de compilación de la aplicación.
  \item El firmware define \texttt{-DARDUINO\_USB\_CDC\_ON\_BOOT=1}. Con ese
        \textit{flag}, el objeto \texttt{Serial} de Arduino se enruta al
        USB nativo en vez del UART0 físico (GPIO43/44). Si la PCB propia usa
        un adaptador UART externo y no tiene el USB nativo conectado a
        nada, ese \textit{flag} debe quitarse —de lo contrario el monitor
        serie del adaptador UART no mostrará nada, aunque el firmware esté
        corriendo correctamente.
\end{enumerate}

\subsection{Circuito mínimo para programar/arrancar}
Cualquier placa propia con el ESP32-S3 necesita, como mínimo: \texttt{EN}
(reset) con \textit{pull-up} a 3.3\,V y un capacitor de $\sim$1\,$\mu$F a
GND; \texttt{GPIO0} (BOOT) con \textit{pull-up} a 3.3\,V y un botón a GND
para forzar el modo \textit{bootloader} manualmente; y, opcionalmente, un
circuito de auto-reset (2 transistores desde DTR/RTS del adaptador UART
hacia EN y GPIO0) para que PlatformIO entre solo al modo de programación sin
botones.

% ============================================================
\section{Solución de problemas comunes}
% ============================================================

\begin{itemize}
  \item \textbf{La subida se queda en "Connecting..." indefinidamente.}
        Entrar manualmente en modo \textit{bootloader} (mantener BOOT,
        pulsar y soltar RESET/EN, soltar BOOT, y recién ahí lanzar la
        subida). Verificar también que ningún otro programa tenga el
        puerto COM abierto.
  \item \textbf{El monitor serie no muestra nada tras programar.} Revisar
        que \texttt{ARDUINO\_USB\_CDC\_ON\_BOOT} coincida con el hardware
        real (Sección~\ref{sec:build}).
  \item \textbf{\texttt{sdSelectCard(): Select Failed} / \texttt{SD=0} en el
        arranque.} Casi siempre es cableado/alimentación: revisar tarjeta
        insertada, tensión correcta del módulo lector, continuidad de los 4
        pines SPI + GND, formato FAT32, y el \textit{pull-up} de
        10\,k$\Omega$ en CS si se soldó directo sin módulo.
  \item \textbf{El panel web no carga en el celular.} Escribir
        \texttt{http://192.168.4.1} con el esquema explícito —varios
        navegadores intentan \texttt{https://} por defecto y el dispositivo
        no tiene TLS—. Algunos celulares además desconectan una red WiFi
        sin salida a internet a menos que se les indique quedarse
        conectados igual.
\end{itemize}
